\documentclass[a4paper,10pt]{article}
\usepackage[margin=1.5cm]{geometry}
\usepackage{multicol}
\usepackage{hyperref}
\usepackage{enumitem}
\usepackage{listings}
\usepackage{graphicx}
\usepackage{longtable}

\setlength{\parindent}{0pt}

\begin{document}


\begin{center}
    {\Large \texttt{Comandos y herramientas Frecuentes en Troubleshooting y sus usos}}\\
\end{center}

\tableofcontents
\newpage

\section{Introducción}

En este documento se encuentran documentados algunos de los comandos y herramientas
más utilizados durante el proceso de \textbf{IT Troubleshooting} en sistemas
Linux, Windows y macOS.

El objetivo principal no es describir exhaustivamente cada comando, sino servir
como una guía de referencia rápida durante la resolución de incidencias.
Cada herramienta se presenta indicando su función principal, sintaxis básica,
argumentos de mayor utilidad, particularidades y el sistema operativo donde
está disponible.

Los comandos se agrupan siguiendo el flujo habitual de un proceso de
troubleshooting, comenzando por la recolección de información, continuando con
la verificación de conectividad, análisis de red, revisión de servicios,
diagnóstico de protocolos como HTTP, HTTPS y TLS, revisión de registros del
sistema y finalizando con herramientas de monitoreo y recuperación.

Es perfectamente válido utilizar comandos compuestos o combinados mediante
pipes, redirecciones o herramientas auxiliares cuando éstos faciliten el
diagnóstico del problema.

\vspace{0.5cm}


%%%%%%%%%%%%%%%%%%%%%%%%%%%%%%%%%%%%%%%%%%%%%%%%%%%%%%%%%%%%%%%%%%%%%%%%%%%%%%%

\section{Listado de comandos y herramientas y sus usos}

El siguiente listado está organizado conforme al flujo habitual de un proceso
de troubleshooting.

Cada subsección agrupa herramientas que permiten obtener cierto tipo de
información o realizar una acción específica para aislar la causa del problema.

En cada tabla se incluyen:

\begin{itemize}
\item El comando o herramienta.
\item Su sintaxis más común.
\item Los argumentos más utilizados durante troubleshooting.
\item Alguna característica importante.
\item Los sistemas operativos donde puede utilizarse.
\end{itemize}

%%%%%%%%%%%%%%%%%%%%%%%%%%%%%%%%%%%%%%%%%%%%%%%%%%%%%%%%%%%%%%%%%%%%%%%%%%%%%%%

\subsection{Paso 1. Recolección de información}

Antes de realizar cualquier modificación sobre el sistema, resulta recomendable
obtener la mayor cantidad de información posible acerca del equipo afectado.

En esta etapa se busca conocer:

\begin{itemize}
\item Sistema operativo.
\item Versión.
\item Kernel.
\item Usuario actual.
\item Arquitectura.
\item Hardware.
\item Memoria disponible.
\item Espacio en disco.
\item Variables de entorno.
\item Hora del sistema.
\end{itemize}

Esta información permite identificar rápidamente incompatibilidades,
restricciones de permisos, problemas de capacidad o configuraciones
incorrectas.

%%%%%%%%%%%%%%%%%%%%%%%%%%%%%%%%%%%%%%%%%%%%%%%%%%%%%%%%%%%%%%%%%%%%%%%%%%%%%%%

\subsubsection{Información general del sistema}

Los siguientes comandos permiten conocer las características generales del
equipo.

\begin{longtable}{|p{2.8cm}|p{3.1cm}|p{4.3cm}|p{5.5cm}|p{1.5cm}|}

\hline
\textbf{Comando} &
\textbf{Sintaxis} &
\textbf{Argumentos útiles} &
\textbf{Particularidad} &
\textbf{OS}
\\
\hline
\endhead

hostname &
hostname &
-f, -I &
Obtiene el nombre del equipo y, dependiendo del sistema, las direcciones IP configuradas. &
Linux/macOS
\\
\hline

hostname &
hostname &
N/A &
Obtiene el nombre del equipo. &
Windows
\\
\hline

uname &
uname -a &
-r, -m, -n &
Muestra versión del kernel, arquitectura y otra información del sistema. &
Linux/macOS
\\
\hline

systeminfo &
systeminfo &
N/A &
Proporciona información completa del sistema operativo, memoria, BIOS, dominio, hotfixes instalados y hardware. &
Windows
\\
\hline

sw\_vers &
sw\_vers &
-productVersion &
Obtiene la versión de macOS instalada. &
macOS
\\
\hline

hostnamectl &
hostnamectl &
status &
Muestra información del hostname, sistema operativo, kernel y arquitectura. &
Linux
\\
\hline

lsb\_release &
lsb\_release -a &
-d &
Obtiene información de la distribución Linux instalada. &
Linux
\\
\hline

ver &
ver &
N/A &
Muestra la versión de Windows desde consola. &
Windows
\\
\hline

winver &
winver &
N/A &
Abre la ventana con la versión gráfica de Windows. &
Windows
\\
\hline

\end{longtable}

%%%%%%%%%%%%%%%%%%%%%%%%%%%%%%%%%%%%%%%%%%%%%%%%%%%%%%%%%%%%%%%%%%%%%%%%%%%%%%%

\subsubsection{Usuario actual y permisos}

Es importante conocer con qué usuario se está trabajando y si éste posee los
privilegios necesarios para realizar cambios en el sistema.

\begin{longtable}{|p{2.8cm}|p{3.1cm}|p{4.3cm}|p{5.5cm}|p{1.5cm}|}

\hline
Comando &
Sintaxis &
Argumentos útiles &
Particularidad &
OS
\\
\hline
\endhead

whoami &
whoami &
N/A &
Muestra el usuario actual. &
Todos
\\
\hline

id &
id &
id usuario &
Obtiene UID, GID y grupos a los que pertenece un usuario. Muy útil para problemas de permisos. &
Linux
\\
\hline

groups &
groups &
usuario &
Lista los grupos de un usuario. &
Linux/macOS
\\
\hline

who &
who &
-a &
Lista los usuarios conectados al sistema. &
Linux/macOS
\\
\hline

w &
w &
-h &
Muestra usuarios conectados y procesos que ejecutan. &
Linux
\\
\hline

query user &
query user &
N/A &
Lista las sesiones activas en Windows. &
Windows
\\
\hline

logname &
logname &
N/A &
Obtiene el usuario que inició la sesión. &
Linux
\\
\hline

sudo -l &
sudo -l &
-U usuario &
Permite conocer los privilegios sudo disponibles. &
Linux
\\
\hline

net user &
net user &
nombreUsuario &
Obtiene información sobre usuarios locales de Windows. &
Windows
\\
\hline

whoami /groups &
whoami /groups &
N/A &
Muestra los grupos de seguridad del usuario actual. &
Windows
\\
\hline

whoami /priv &
whoami /priv &
N/A &
Lista los privilegios asignados al usuario actual. &
Windows
\\
\hline

\end{longtable}

%%%%%%%%%%%%%%%%%%%%%%%%%%%%%%%%%%%%%%%%%%%%%%%%%%%%%%%%%%%%%%%%%%%%%%%%%%%%%%%

\subsubsection{Fecha y hora}

Una configuración incorrecta de fecha u hora puede ocasionar problemas de
autenticación, certificados TLS, Kerberos, sincronización de archivos y
servicios distribuidos.

\begin{longtable}{|p{2.8cm}|p{3.1cm}|p{4.3cm}|p{5.5cm}|p{1.5cm}|}

\hline
Comando &
Sintaxis &
Argumentos útiles &
Particularidad &
OS
\\
\hline
\endhead

date &
date &
-u &
Obtiene la fecha y hora del sistema. &
Linux/macOS
\\
\hline

timedatectl &
timedatectl &
status &
Muestra la configuración de fecha, zona horaria y estado de sincronización NTP. &
Linux
\\
\hline

Get-Date &
Get-Date &
-Format &
Obtiene la fecha desde PowerShell. &
Windows
\\
\hline

tzutil &
tzutil /g &
/s &
Consulta o modifica la zona horaria. &
Windows
\\
\hline

w32tm &
w32tm /query /status &
/resync &
Consulta el estado del servicio de sincronización NTP de Windows. &
Windows
\\
\hline

systemsetup &
systemsetup -gettimezone &
-settimezone &
Consulta o modifica la zona horaria en macOS. &
macOS
\\
\hline

\end{longtable}

%%%%%%%%%%%%%%%%%%%%%%%%%%%%%%%%%%%%%%%%%%%%%%%%%%%%%%%%%%%%%%%%%%%%%%%%%%%%%%%

\subsubsection{Información del hardware}

Estas herramientas permiten conocer las características físicas del equipo.

\begin{longtable}{|p{2.8cm}|p{3.1cm}|p{4.3cm}|p{5.5cm}|p{1.5cm}|}

\hline
Comando &
Sintaxis &
Argumentos útiles &
Particularidad &
OS
\\
\hline
\endhead

lscpu &
lscpu &
-e &
Obtiene información detallada del procesador. &
Linux
\\
\hline

cat /proc/cpuinfo &
cat /proc/cpuinfo &
grep model &
Información detallada del CPU. &
Linux
\\
\hline

lsblk &
lsblk &
-f, -o &
Lista los dispositivos de almacenamiento conectados. &
Linux
\\
\hline

lshw &
sudo lshw &
-short &
Inventario completo del hardware instalado. &
Linux
\\
\hline

dmidecode &
sudo dmidecode &
-t memory, -t bios &
Obtiene información directamente del BIOS/UEFI. &
Linux
\\
\hline

wmic cpu &
wmic cpu get name &
N/A &
Información del procesador. &
Windows
\\
\hline

wmic memorychip &
wmic memorychip &
N/A &
Obtiene información de la memoria RAM instalada. &
Windows
\\
\hline

system\_profiler &
system\_profiler SPHardwareDataType &
SPMemoryDataType &
Inventario completo del hardware en macOS. &
macOS
\\
\hline



sysctl &
sysctl hw &
-a &
Consulta parámetros del kernel y cierta información del hardware. Muy útil para verificar configuraciones del sistema y límites del kernel. &
macOS/Linux
\\
\hline

lspci &
lspci &
-v, -vv, -k &
Lista los dispositivos conectados al bus PCI (adaptadores de red, GPU, controladores SATA, etc.). Muy útil para verificar que el hardware sea detectado correctamente. &
Linux
\\
\hline

lsusb &
lsusb &
-v &
Lista los dispositivos USB conectados al equipo. Útil cuando un periférico no es reconocido. &
Linux
\\
\hline

diskutil &
diskutil list &
info &
Permite listar discos y particiones disponibles en macOS. &
macOS
\\
\hline

Get-ComputerInfo &
Get-ComputerInfo &
-property * &
Obtiene un inventario muy completo del sistema mediante PowerShell. &
Windows
\\
\hline

\end{longtable}

%%%%%%%%%%%%%%%%%%%%%%%%%%%%%%%%%%%%%%%%%%%%%%%%%%%%%%%%%%%%%%%%%%%%%%%%%%%%%%%

\subsubsection{Uso de memoria}

Problemas relacionados con falta de memoria RAM o memoria virtual pueden provocar
lentitud, terminación inesperada de procesos o fallas al iniciar servicios.

Los siguientes comandos permiten conocer el estado actual de la memoria.

\begin{longtable}{|p{2.8cm}|p{3cm}|p{4cm}|p{5.6cm}|p{1.5cm}|}

\hline
Comando &
Sintaxis &
Argumentos útiles &
Particularidad &
OS
\\
\hline
\endhead

free &
free -h &
-m, -g &
Muestra la memoria RAM y Swap utilizada. El argumento \texttt{-h} facilita su lectura mostrando unidades apropiadas. &
Linux
\\
\hline

vmstat &
vmstat &
1, -s &
Muestra estadísticas de memoria, CPU y procesos. Muy útil para observar el comportamiento del sistema durante algunos segundos. &
Linux/macOS
\\
\hline

top &
top &
-d, -o &
Permite observar el consumo de memoria de los procesos en tiempo real. &
Linux/macOS
\\
\hline

htop &
htop &
N/A &
Versión interactiva de \texttt{top}. Facilita ordenar procesos por uso de memoria o CPU. &
Linux
\\
\hline

cat /proc/meminfo &
cat /proc/meminfo &
grep Mem &
Información detallada de la memoria administrada por el kernel. &
Linux
\\
\hline

vm\_stat &
vm\_stat &
N/A &
Muestra estadísticas detalladas de memoria virtual en macOS. &
macOS
\\
\hline

Activity Monitor &
N/A &
Memory &
Herramienta gráfica para revisar el consumo de memoria en macOS. &
macOS
\\
\hline

taskmgr &
taskmgr &
N/A &
Administrador de tareas de Windows. Permite revisar el consumo de memoria por proceso. &
Windows
\\
\hline

Get-Process &
Get-Process &
Sort WorkingSet &
Permite ordenar procesos por consumo de memoria desde PowerShell. &
Windows
\\
\hline

\end{longtable}

%%%%%%%%%%%%%%%%%%%%%%%%%%%%%%%%%%%%%%%%%%%%%%%%%%%%%%%%%%%%%%%%%%%%%%%%%%%%%%%

\subsubsection{Uso de almacenamiento}

La falta de espacio en disco suele ocasionar errores de aplicaciones,
servicios que dejan de funcionar, imposibilidad de generar logs,
errores de bases de datos y fallas durante actualizaciones.

\begin{longtable}{|p{2.8cm}|p{3cm}|p{4cm}|p{5.6cm}|p{1.5cm}|}

\hline
Comando &
Sintaxis &
Argumentos útiles &
Particularidad &
OS
\\
\hline
\endhead

df &
df -h &
-i &
Muestra el espacio libre por sistema de archivos. El argumento \texttt{-i} permite revisar el consumo de inodos. &
Linux/macOS
\\
\hline

du &
du -sh directorio &
-h, --max-depth &
Calcula el tamaño ocupado por archivos o directorios. Muy útil para localizar qué consume espacio. &
Linux/macOS
\\
\hline

ncdu &
ncdu &
N/A &
Herramienta interactiva para localizar directorios con mayor consumo de almacenamiento. &
Linux
\\
\hline

lsblk &
lsblk -f &
-o &
Lista discos, particiones y sistemas de archivos disponibles. &
Linux
\\
\hline

mount &
mount &
N/A &
Lista los sistemas de archivos montados. &
Linux/macOS
\\
\hline

findmnt &
findmnt &
N/A &
Presenta de forma estructurada la información de los puntos de montaje. &
Linux
\\
\hline

diskutil &
diskutil list &
info &
Información sobre discos y particiones en macOS. &
macOS
\\
\hline

wmic logicaldisk &
wmic logicaldisk get size,freespace,caption &
N/A &
Consulta el espacio disponible de las unidades en Windows. &
Windows
\\
\hline

Get-PSDrive &
Get-PSDrive &
N/A &
Lista las unidades disponibles y el espacio libre mediante PowerShell. &
Windows
\\
\hline

\end{longtable}

%%%%%%%%%%%%%%%%%%%%%%%%%%%%%%%%%%%%%%%%%%%%%%%%%%%%%%%%%%%%%%%%%%%%%%%%%%%%%%%

\subsubsection{Variables de entorno}

Muchas aplicaciones dependen de variables de entorno para localizar
bibliotecas, archivos de configuración, certificados o ejecutables.

Una configuración incorrecta suele provocar errores aparentemente
inexplicables.

\begin{longtable}{|p{2.8cm}|p{3cm}|p{4cm}|p{5.6cm}|p{1.5cm}|}

\hline
Comando &
Sintaxis &
Argumentos útiles &
Particularidad &
OS
\\
\hline
\endhead

env &
env &
grep VARIABLE &
Lista todas las variables de entorno del proceso actual. &
Linux/macOS
\\
\hline

printenv &
printenv &
PATH &
Permite consultar una variable específica. &
Linux/macOS
\\
\hline

echo &
echo \$PATH &
\$HOME &
Consulta el contenido de una variable de entorno. &
Linux/macOS
\\
\hline

set &
set &
N/A &
Lista variables de entorno en CMD. &
Windows
\\
\hline

Get-ChildItem Env: &
Get-ChildItem Env: &
N/A &
Lista todas las variables de entorno desde PowerShell. &
Windows
\\
\hline

echo &
echo \%PATH\% &
\%TEMP\% &
Consulta variables específicas desde CMD. &
Windows
\\
\hline

which &
which programa &
-a &
Indica qué ejecutable será utilizado al invocar un comando. Muy útil para detectar conflictos entre múltiples instalaciones. &
Linux/macOS
\\
\hline

whereis &
whereis programa &
N/A &
Localiza binarios, código fuente y páginas de manual. &
Linux
\\
\hline

where &
where programa &
N/A &
Busca la ubicación de un ejecutable en Windows. &
Windows
\\
\hline

type &
type comando &
-a &
Indica si un comando corresponde a un alias, función o ejecutable. Muy útil durante troubleshooting de shells. &
Linux
\\
\hline

\end{longtable}

%%%%%%%%%%%%%%%%%%%%%%%%%%%%%%%%%%%%%%%%%%%%%%%%%%%%%%%%%%%%%%%%%%%%%%%%%%%%%%%

\subsubsection{Información de procesos}

Antes de comenzar a revisar conectividad o servicios, resulta recomendable
verificar si el proceso o aplicación realmente se encuentra en ejecución.

En muchas ocasiones el problema no corresponde a la red, sino a que el
servicio nunca inició correctamente.

\begin{longtable}{|p{2.8cm}|p{3cm}|p{4cm}|p{5.6cm}|p{1.5cm}|}

\hline
Comando &
Sintaxis &
Argumentos útiles &
Particularidad &
OS
\\
\hline
\endhead

ps &
ps aux &
-ef &
Lista los procesos activos del sistema. Constituye una de las herramientas fundamentales de troubleshooting en Linux. &
Linux/macOS
\\
\hline

pgrep &
pgrep proceso &
-a, -f &
Busca procesos por nombre o por línea completa de ejecución. &
Linux
\\
\hline

pidof &
pidof nginx &
N/A &
Obtiene rápidamente el PID de un proceso. &
Linux
\\
\hline

pstree &
pstree &
-p &
Representa gráficamente la jerarquía de procesos. Muy útil para identificar procesos padre e hijos. &
Linux
\\
\hline

top &
top &
-o &
Monitor interactivo de procesos en tiempo real. &
Linux/macOS
\\
\hline

htop &
htop &
N/A &
Versión mejorada e interactiva de top. &
Linux
\\
\hline

tasklist &
tasklist &
/svc, /FI &
Lista los procesos activos en Windows. Puede filtrar procesos por nombre o servicio asociado. &
Windows
\\
\hline

Get-Process &
Get-Process &
-Name &
Obtiene procesos mediante PowerShell. &
Windows
\\
\hline

Get-Service &
Get-Service &
-Name, -DisplayName &
Permite identificar rápidamente si un proceso está asociado a un servicio del sistema. Muy útil para determinar si un servicio se encuentra detenido o en ejecución. &
Windows
\\
\hline

procmon &
procmon &
Filtros &
Herramienta de Sysinternals que permite monitorear en tiempo real procesos, acceso a archivos, registro y red. Es una de las herramientas más utilizadas durante troubleshooting avanzado en Windows. &
Windows
\\
\hline

Process Explorer &
procexp &
N/A &
Herramienta de Sysinternals que proporciona información detallada sobre procesos, DLL cargadas, handles abiertos y uso de recursos. &
Windows
\\
\hline

Activity Monitor &
Activity Monitor &
CPU, Memory &
Administrador gráfico de procesos en macOS. Permite revisar consumo de CPU, memoria, disco y red. &
macOS
\\
\hline

kill &
kill PID &
-9, -15, -HUP &
Envía señales a un proceso. Durante troubleshooting suele utilizarse para finalizar procesos bloqueados o solicitar que recarguen su configuración. &
Linux/macOS
\\
\hline

pkill &
pkill nombreProceso &
-f &
Finaliza procesos buscando por nombre. Evita tener que buscar previamente el PID. &
Linux
\\
\hline

killall &
killall nombreProceso &
-v &
Finaliza todos los procesos con el mismo nombre. Disponible en Linux y macOS. &
Linux/macOS
\\
\hline

taskkill &
taskkill /PID PID &
/F, /IM &
Finaliza procesos en Windows mediante CMD. Puede utilizarse indicando el PID o el nombre del ejecutable. &
Windows
\\
\hline

Stop-Process &
Stop-Process -Id PID &
-Force &
Equivalente en PowerShell para finalizar procesos. &
Windows
\\
\hline

\end{longtable}

%%%%%%%%%%%%%%%%%%%%%%%%%%%%%%%%%%%%%%%%%%%%%%%%%%%%%%%%%%%%%%%%%%%%%%%%%%%%%%%

\subsubsection{Información de servicios}

En muchos escenarios el problema no se encuentra en la aplicación en sí, sino en
que el servicio correspondiente nunca inició, se encuentra detenido o falló
durante el arranque del sistema.

Las siguientes herramientas permiten consultar el estado de los servicios,
iniciarlos, detenerlos o reiniciarlos.

\begin{longtable}{|p{2.8cm}|p{3cm}|p{4cm}|p{5.6cm}|p{1.5cm}|}

\hline
Comando &
Sintaxis &
Argumentos útiles &
Particularidad &
OS
\\
\hline
\endhead

systemctl &
systemctl status servicio &
start, stop, restart, enable, disable, is-active &
Principal herramienta para administrar servicios en sistemas Linux basados en systemd. &
Linux
\\
\hline

service &
service nombreServicio status &
start, stop, restart &
Interfaz compatible con sistemas Linux más antiguos o distribuciones que mantienen scripts SysV. &
Linux
\\
\hline

journalctl &
journalctl -u servicio &
-f, -b, --since &
Permite consultar exclusivamente los registros asociados a un servicio administrado por systemd. Muy útil cuando un servicio falla al iniciar. &
Linux
\\
\hline

systemctl list-units &
systemctl list-units --type=service &
--failed &
Lista los servicios cargados y permite identificar rápidamente aquellos que fallaron durante el arranque. &
Linux
\\
\hline

launchctl &
launchctl list &
load, unload, kickstart &
Administrador de servicios y daemons en macOS. &
macOS
\\
\hline

sc &
sc query servicio &
start, stop, qc, config &
Administrador de servicios desde la línea de comandos de Windows. Permite consultar configuración y estado. &
Windows
\\
\hline

net start &
net start &
nombreServicio &
Lista los servicios actualmente iniciados o inicia uno específico. &
Windows
\\
\hline

net stop &
net stop servicio &
N/A &
Detiene un servicio en Windows. &
Windows
\\
\hline

Get-Service &
Get-Service &
Status, Name &
Obtiene el estado de los servicios mediante PowerShell. Permite filtrar fácilmente por nombre o estado. &
Windows
\\
\hline

Restart-Service &
Restart-Service servicio &
-Force &
Reinicia un servicio desde PowerShell. Muy útil durante troubleshooting remoto. &
Windows
\\
\hline

services.msc &
services.msc &
N/A &
Consola gráfica para administrar servicios en Windows. Facilita revisar dependencias y tipo de inicio. &
Windows
\\
\hline

\end{longtable}

%%%%%%%%%%%%%%%%%%%%%%%%%%%%%%%%%%%%%%%%%%%%%%%%%%%%%%%%%%%%%%%%%%%%%%%%%%%%%%%

\subsubsection{Información del sistema de archivos}

Una gran cantidad de incidencias provienen de problemas relacionados con
permisos, archivos faltantes, puntos de montaje incorrectos o configuraciones
ubicadas en rutas inesperadas.

Los siguientes comandos permiten explorar el sistema de archivos.

\begin{longtable}{|p{2.8cm}|p{3cm}|p{4cm}|p{5.6cm}|p{1.5cm}|}

\hline
Comando &
Sintaxis &
Argumentos útiles &
Particularidad &
OS
\\
\hline
\endhead

pwd &
pwd &
N/A &
Muestra el directorio actual de trabajo. Fundamental para evitar operar sobre rutas incorrectas. &
Linux/macOS
\\
\hline

cd &
cd ruta &
.., ~ &
Permite cambiar el directorio de trabajo. &
Todos
\\
\hline

ls &
ls -lah &
-R, -ltr &
Lista el contenido de un directorio. Los argumentos permiten visualizar permisos, tamaños y ordenar por fecha de modificación. &
Linux/macOS
\\
\hline

dir &
dir &
/A, /O &
Equivalente de \texttt{ls} en Windows CMD. &
Windows
\\
\hline

tree &
tree &
-L, /F &
Representa gráficamente la estructura de directorios. Muy útil para localizar configuraciones o entender la organización de un proyecto. &
Todos
\\
\hline

stat &
stat archivo &
-c &
Obtiene información detallada de un archivo: permisos, propietario, fechas de acceso y modificación, tamaño e inodos. &
Linux/macOS
\\
\hline

attrib &
attrib archivo &
+R, -R &
Consulta o modifica atributos de archivos en Windows. &
Windows
\\
\hline

file &
file archivo &
N/A &
Determina el tipo real de un archivo independientemente de su extensión. Muy útil para identificar archivos corruptos o con extensiones incorrectas. &
Linux/macOS
\\
\hline

realpath &
realpath archivo &
N/A &
Obtiene la ruta absoluta de un archivo o directorio. &
Linux/macOS
\\
\hline

readlink &
readlink -f enlace &
-f &
Resuelve enlaces simbólicos y muestra el archivo real al que apuntan. &
Linux
\\
\hline


find &
find ruta criterios &
-name, -iname, -type, -size, -mtime, -mmin, -user, -perm, -exec &
Herramienta principal para localizar archivos y directorios utilizando múltiples criterios de búsqueda. Fundamental durante tareas de troubleshooting. &
Linux/macOS
\\
\hline

locate &
locate archivo &
-i &
Busca archivos utilizando una base de datos previamente indexada. Mucho más rápido que \texttt{find}, aunque puede no reflejar cambios recientes. &
Linux
\\
\hline

updatedb &
updatedb &
N/A &
Actualiza la base de datos utilizada por \texttt{locate}. &
Linux
\\
\hline

whereis &
whereis comando &
-b, -m &
Localiza el binario, código fuente y páginas del manual asociadas a un programa. &
Linux
\\
\hline

which &
which programa &
-a &
Indica qué ejecutable será utilizado al invocar un comando. Muy útil para detectar múltiples instalaciones de un mismo programa. &
Linux/macOS
\\
\hline

where &
where programa &
/R &
Busca ejecutables dentro del PATH o dentro de un directorio específico. &
Windows
\\
\hline

findstr &
findstr patrón archivo &
/S, /I, /N &
Busca cadenas de texto dentro de archivos. Equivalente básico de \texttt{grep} en Windows. &
Windows
\\
\hline

grep &
grep patrón archivo &
-r, -i, -n, -v, -E, -A, -B, -C &
Permite buscar texto dentro de archivos. Es una de las herramientas más utilizadas durante troubleshooting para localizar mensajes de error o configuraciones. &
Linux/macOS
\\
\hline

egrep &
egrep expresión archivo &
N/A &
Versión basada en expresiones regulares extendidas. Actualmente suele reemplazarse por \texttt{grep -E}. &
Linux/macOS
\\
\hline

zgrep &
zgrep patrón archivo.gz &
-i &
Permite buscar texto directamente dentro de archivos comprimidos con gzip sin necesidad de descomprimirlos. &
Linux
\\
\hline

strings &
strings archivo &
-a &
Extrae cadenas ASCII imprimibles de archivos binarios. Muy útil para identificar rutas, URLs, mensajes o configuraciones embebidas. &
Linux/macOS
\\
\hline

head &
head archivo &
-n &
Muestra las primeras líneas de un archivo. &
Linux/macOS
\\
\hline

tail &
tail archivo &
-f, -n &
Muestra las últimas líneas de un archivo. El argumento \texttt{-f} permite seguir un archivo en tiempo real. Muy utilizado para monitorear logs. &
Linux/macOS
\\
\hline

more &
more archivo &
N/A &
Visualiza archivos página por página. &
Todos
\\
\hline

less &
less archivo &
-N &
Permite navegar archivos grandes de manera interactiva. Constituye una de las herramientas más utilizadas para revisar logs. &
Linux/macOS
\\
\hline

cat &
cat archivo &
-n &
Muestra el contenido completo de uno o varios archivos. &
Linux/macOS
\\
\hline

type &
type archivo &
N/A &
Muestra el contenido de un archivo desde CMD. &
Windows
\\
\hline

sort &
sort archivo &
-r, -u &
Ordena el contenido de un archivo. Puede eliminar líneas duplicadas utilizando \texttt{-u}. &
Linux/macOS
\\
\hline

uniq &
uniq archivo &
-c, -d &
Elimina líneas consecutivas duplicadas o cuenta ocurrencias. Suele utilizarse junto con \texttt{sort}. &
Linux/macOS
\\
\hline

wc &
wc archivo &
-l, -w, -c &
Cuenta líneas, palabras y caracteres. Muy útil para validar el tamaño de un archivo o el número de resultados obtenidos. &
Linux/macOS
\\
\hline

cut &
cut &
-d, -f &
Extrae columnas específicas de un archivo delimitado por caracteres. &
Linux/macOS
\\
\hline

paste &
paste &
-d &
Une líneas de varios archivos en columnas. &
Linux/macOS
\\
\hline

tr &
tr &
-d, -s &
Reemplaza o elimina caracteres de la entrada estándar. &
Linux/macOS
\\
\hline

sed &
sed &
-n, -i, -e &
Editor de flujo ampliamente utilizado para modificar archivos de configuración o transformar texto durante el troubleshooting. &
Linux/macOS
\\
\hline

awk &
awk &
-F &
Lenguaje de procesamiento de texto orientado a columnas. Muy útil para analizar logs, archivos CSV y salidas de comandos. &
Linux/macOS
\\
\hline

xargs &
xargs &
-I, -0 &
Permite construir comandos utilizando la salida de otro comando. Muy utilizado junto con \texttt{find}. &
Linux/macOS
\\
\hline

\end{longtable}

%%%%%%%%%%%%%%%%%%%%%%%%%%%%%%%%%%%%%%%%%%%%%%%%%%%%%%%%%%%%%%%%%%%%%%%%%%%%%%%

\subsubsection{Permisos y propietarios}

Problemas relacionados con permisos insuficientes constituyen una de las causas
más frecuentes de incidencias en sistemas operativos. Los siguientes comandos
permiten consultar o modificar permisos, propietarios y listas de control de acceso.

\begin{longtable}{|p{2.8cm}|p{3cm}|p{4cm}|p{5.6cm}|p{1.5cm}|}

\hline
Comando &
Sintaxis &
Argumentos útiles &
Particularidad &
OS
\\
\hline
\endhead

ls &
ls -l &
-a, -h &
Permite visualizar permisos, propietario y grupo de archivos y directorios. &
Linux/macOS
\\
\hline

stat &
stat archivo &
-c &
Obtiene información detallada sobre permisos, propietario, fechas e inodos. &
Linux/macOS
\\
\hline

chmod &
chmod modo archivo &
-R &
Modifica los permisos de archivos y directorios utilizando notación simbólica u octal. &
Linux/macOS
\\
\hline

chown &
chown usuario archivo &
-R &
Modifica el propietario y grupo de un archivo o directorio. &
Linux/macOS
\\
\hline

chgrp &
chgrp grupo archivo &
-R &
Modifica únicamente el grupo propietario. &
Linux/macOS
\\
\hline

umask &
umask &
-S &
Consulta o modifica la máscara de permisos predeterminada para nuevos archivos. &
Linux/macOS
\\
\hline

getfacl &
getfacl archivo &
-R &
Consulta listas de control de acceso (ACL). Muy útil cuando los permisos tradicionales no explican el comportamiento observado. &
Linux
\\
\hline

setfacl &
setfacl &
-m, -x, -R &
Permite agregar, modificar o eliminar entradas ACL. &
Linux
\\
\hline

icacls &
icacls archivo &
/grant, /remove, /inheritance &
Consulta y modifica ACL en Windows. &
Windows
\\
\hline

takeown &
takeown /F archivo &
/R &
Permite tomar posesión de archivos o directorios cuando el usuario carece de permisos suficientes. &
Windows
\\
\hline

cacls &
cacls archivo &
/E &
Herramienta heredada para administrar permisos NTFS. Actualmente se recomienda utilizar \texttt{icacls}. &
Windows
\\
\hline

Get-Acl &
Get-Acl ruta &
N/A &
Obtiene las listas de control de acceso desde PowerShell. &
Windows
\\
\hline

Set-Acl &
Set-Acl &
N/A &
Modifica permisos mediante PowerShell. &
Windows
\\
\hline

\end{longtable}

%%%%%%%%%%%%%%%%%%%%%%%%%%%%%%%%%%%%%%%%%%%%%%%%%%%%%%%%%%%%%%%%%%%%%%%%%%%%%%%

\subsection{Paso 2. Verificación del estado de la red}

Una vez recopilada la información básica del sistema, el siguiente paso consiste
en verificar el estado de la configuración de red.

Antes de comprobar la conectividad hacia un servicio remoto es recomendable
confirmar que la interfaz de red se encuentre habilitada, que posea una dirección
IP válida, una máscara de red correcta, un gateway configurado y los servidores
DNS esperados.

Los comandos de esta sección permiten inspeccionar la configuración local del
equipo antes de comenzar a diagnosticar problemas de comunicación.

%%%%%%%%%%%%%%%%%%%%%%%%%%%%%%%%%%%%%%%%%%%%%%%%%%%%%%%%%%%%%%%%%%%%%%%%%%%%%%%

\subsubsection{Interfaces de red}

Las siguientes herramientas permiten consultar las interfaces disponibles,
su estado, direcciones IP asignadas y otra información relevante de la configuración
de red.
%%%%%%%%%%%%%%%%%%

\begin{longtable}{|p{2.8cm}|p{3cm}|p{4cm}|p{5.6cm}|p{1.5cm}|}

\hline
Comando &
Sintaxis &
Argumentos útiles &
Particularidad &
OS
\\
\hline
\endhead

ip &
ip addr &
link, addr, -br &
Herramienta moderna para consultar y administrar interfaces de red, direcciones IP y enlaces. Sustituye a \texttt{ifconfig} en la mayoría de las distribuciones Linux. &
Linux
\\
\hline

ip link &
ip link show &
set up, set down &
Consulta el estado de las interfaces y permite habilitarlas o deshabilitarlas. &
Linux
\\
\hline

ip addr &
ip addr show &
add, del &
Permite consultar, agregar o eliminar direcciones IP asignadas a una interfaz. &
Linux
\\
\hline

ifconfig &
ifconfig &
-a &
Herramienta clásica para visualizar y configurar interfaces de red. Aunque está obsoleta en muchas distribuciones, sigue presente en numerosos sistemas. &
Linux/macOS
\\
\hline

ipconfig &
ipconfig &
/all, /release, /renew, /flushdns, /displaydns &
Muestra la configuración de red del equipo y permite renovar direcciones obtenidas mediante DHCP. &
Windows
\\
\hline

networksetup &
networksetup -listallhardwareports &
-getinfo, -setdnsservers &
Permite consultar y modificar la configuración de red en macOS. &
macOS
\\
\hline

ifstatus &
ifstatus interfaz &
N/A &
Consulta el estado de una interfaz en algunas distribuciones Linux. &
Linux
\\
\hline

nmcli &
nmcli device status &
connection show, device show &
Interfaz de línea de comandos de NetworkManager. Muy utilizada en distribuciones modernas. &
Linux
\\
\hline

hostname -I &
hostname -I &
N/A &
Muestra rápidamente las direcciones IPv4 e IPv6 asignadas al equipo. &
Linux
\\
\hline

Get-NetAdapter &
Get-NetAdapter &
Name, Status &
Lista los adaptadores de red instalados mediante PowerShell. &
Windows
\\
\hline

Get-NetIPAddress &
Get-NetIPAddress &
InterfaceAlias &
Obtiene todas las direcciones IP configuradas en Windows mediante PowerShell. &
Windows
\\
\hline

netsh interface show interface &
netsh interface show interface &
N/A &
Lista las interfaces de red y su estado administrativo y operativo. &
Windows
\\
\hline

\end{longtable}

%%%%%%%%%%%%%%%%%%%%%%%%%%%%%%%%%%%%%%%%%%%%%%%%%%%%%%%%%%%%%%%%%%%%%%%%%%%%%%%

\subsubsection{Configuración IP}

Una configuración IP incorrecta es una de las causas más comunes de problemas de
conectividad. Los siguientes comandos permiten consultar direcciones IPv4,
IPv6, máscaras de red, gateways y servidores DNS configurados.

\begin{longtable}{|p{2.8cm}|p{3cm}|p{4cm}|p{5.6cm}|p{1.5cm}|}

\hline
Comando &
Sintaxis &
Argumentos útiles &
Particularidad &
OS
\\
\hline
\endhead

ip addr show &
ip addr show &
dev interfaz &
Muestra las direcciones IP configuradas en una interfaz específica. &
Linux
\\
\hline

ip -br addr &
ip -br addr &
N/A &
Presenta un resumen compacto de las interfaces y sus direcciones IP. Muy útil durante sesiones remotas. &
Linux
\\
\hline

ifconfig &
ifconfig interfaz &
N/A &
Consulta la dirección IP y parámetros de una interfaz específica. &
Linux/macOS
\\
\hline

ipconfig /all &
ipconfig /all &
N/A &
Presenta la configuración completa de todos los adaptadores de red en Windows. &
Windows
\\
\hline

Get-NetIPAddress &
Get-NetIPAddress &
AddressFamily &
Obtiene información detallada de las direcciones IPv4 e IPv6 configuradas. &
Windows
\\
\hline

hostname -I &
hostname -I &
N/A &
Muestra únicamente las direcciones IP configuradas. &
Linux
\\
\hline

networksetup -getinfo &
networksetup -getinfo Ethernet &
Wi-Fi &
Consulta la configuración IP de una interfaz en macOS. &
macOS
\\
\hline


\end{longtable}
%%%%%%%%

%%%%%%%%%%%%%%%%%%%%%%%%%%%%%%%%%%%%%%%%%%%%%%%%%%%%%%%%%%%%%%%%%%%%%%%%%%%%%%%

\subsubsection{Gateway y tabla de rutas}

Si la configuración IP parece correcta pero no existe conectividad hacia otras
redes, resulta necesario revisar la puerta de enlace predeterminada y la tabla
de enrutamiento del sistema.

\begin{longtable}{|p{2.8cm}|p{3cm}|p{4cm}|p{5.6cm}|p{1.5cm}|}

\hline
Comando &
Sintaxis &
Argumentos útiles &
Particularidad &
OS
\\
\hline
\endhead

ip route &
ip route &
show, add, del &
Consulta y administra la tabla de enrutamiento del sistema. Constituye la herramienta principal para diagnosticar problemas de routing en Linux. &
Linux
\\
\hline

route &
route -n &
add, del &
Herramienta clásica para consultar y modificar rutas. Aún disponible en numerosos sistemas UNIX. &
Linux/macOS
\\
\hline

netstat &
netstat -rn &
-r &
Permite visualizar la tabla de rutas utilizando la utilidad netstat. &
Linux/macOS
\\
\hline

route print &
route print &
N/A &
Muestra la tabla de enrutamiento de Windows. Incluye rutas IPv4 e IPv6. &
Windows
\\
\hline

Get-NetRoute &
Get-NetRoute &
DestinationPrefix &
Consulta las rutas configuradas mediante PowerShell. &
Windows
\\
\hline

ip route get &
ip route get direcciónIP &
N/A &
Indica qué ruta utilizará el sistema para llegar a un destino específico. Muy útil para validar decisiones de enrutamiento. &
Linux
\\
\hline

route get &
route get direcciónIP &
N/A &
Obtiene la ruta utilizada hacia un destino específico en macOS. &
macOS
\\
\hline

\end{longtable}
%%%%%%%%%5

%%%%%%%%%%%%%%%%%%%%%%%%%%%%%%%%%%%%%%%%%%%%%%%%%%%%%%%%%%%%%%%%%%%%%%%%%%%%%%%

\subsubsection{ARP y vecinos}

Antes de transmitir tráfico dentro de una red local, un equipo necesita conocer
la dirección MAC correspondiente a la dirección IP de destino. Para ello utiliza
el protocolo ARP (IPv4) o Neighbor Discovery (IPv6).

Las siguientes herramientas permiten inspeccionar dichas tablas.

\begin{longtable}{|p{2.8cm}|p{3cm}|p{4cm}|p{5.6cm}|p{1.5cm}|}

\hline
Comando &
Sintaxis &
Argumentos útiles &
Particularidad &
OS
\\
\hline
\endhead

arp &
arp -a &
-d &
Muestra la caché ARP del sistema y permite eliminar entradas específicas. &
Linux/macOS/Windows
\\
\hline

ip neigh &
ip neigh show &
flush &
Consulta la tabla de vecinos utilizada por el kernel para IPv4 e IPv6. Sustituye progresivamente al comando \texttt{arp}. &
Linux
\\
\hline

ip -s neigh &
ip -s neigh &
N/A &
Incluye estadísticas adicionales de cada vecino conocido. &
Linux
\\
\hline

netsh interface ip show neighbors &
netsh interface ip show neighbors &
N/A &
Consulta la tabla ARP desde Windows. &
Windows
\\
\hline

Get-NetNeighbor &
Get-NetNeighbor &
InterfaceAlias &
Obtiene la tabla de vecinos mediante PowerShell. &
Windows
\\
\hline

arp -d &
arp -d direcciónIP &
N/A &
Elimina una entrada específica de la caché ARP para forzar una nueva resolución. &
Linux/macOS/Windows
\\
\hline

ip neigh flush &
ip neigh flush all &
dev &
Vacía completamente la tabla de vecinos. Muy útil para diagnosticar problemas de resolución MAC. &
Linux
\\
\hline

\end{longtable}

%%%%%%%%55

%%%%%%%%%%%%%%%%%%%%%%%%%%%%%%%%%%%%%%%%%%%%%%%%%%%%%%%%%%%%%%%%%%%%%%%%%%%%%%%

\subsection{Paso 3. Verificación de conectividad}

Una vez validada la configuración local de la red, el siguiente paso consiste en
comprobar la conectividad hacia otros dispositivos.

Durante esta etapa se busca responder preguntas como:

\begin{itemize}
\item ¿El equipo tiene comunicación con su gateway?
\item ¿Existe conectividad hacia otros equipos de la misma red?
\item ¿Existe salida hacia Internet?
\item ¿En qué salto se pierde la comunicación?
\item ¿El puerto del servicio realmente está respondiendo?
\item ¿El problema corresponde a la red o únicamente a la aplicación?
\end{itemize}

%%%%%%%%%%%%%%%%%%%%%%%%%%%%%%%%%%%%%%%%%%%%%%%%%%%%%%%%%%%%%%%%%%%%%%%%%%%%%%%

\subsubsection{Verificación básica de conectividad}

Las siguientes herramientas permiten verificar rápidamente si existe
comunicación entre dos equipos utilizando distintos protocolos.

\begin{longtable}{|p{2.8cm}|p{3cm}|p{4cm}|p{5.6cm}|p{1.5cm}|}

\hline
Comando &
Sintaxis &
Argumentos útiles &
Particularidad &
OS
\\
\hline
\endhead

ping &
ping destino &
-c, -i, -s, -W &
Envía solicitudes ICMP Echo Request para verificar conectividad básica, latencia y pérdida de paquetes. Constituye la primera prueba de conectividad en la mayoría de los escenarios. &
Linux/macOS
\\
\hline

ping &
ping destino &
-n, -l, -t, -4, -6 &
Versión disponible en Windows. Permite forzar IPv4 o IPv6 y realizar pruebas continuas. &
Windows
\\
\hline

ping6 &
ping6 destino &
-c &
Permite comprobar conectividad utilizando IPv6. En muchas distribuciones modernas también puede utilizarse mediante \texttt{ping -6}. &
Linux/macOS
\\
\hline

fping &
fping rango &
-a, -g &
Permite realizar pruebas ICMP sobre múltiples equipos simultáneamente. Muy útil durante tareas de inventario o diagnóstico masivo. &
Linux
\\
\hline

arping &
arping direcciónIP &
-c, -I &
Envía solicitudes ARP en lugar de ICMP. Muy útil para comprobar conectividad dentro de la misma red incluso cuando ICMP está bloqueado. &
Linux
\\
\hline


\end{longtable}
%%%%

%%%%%%%%%%%%%%%%%%%%%%%%%%%%%%%%%%%%%%%%%%%%%%%%%%%%%%%%%%%%%%%%%%%%%%%%%%%%%%%

\subsubsection{Análisis de rutas}

Cuando existe pérdida de conectividad, resulta importante determinar en qué
dispositivo o salto se interrumpe la comunicación.

Las siguientes herramientas permiten analizar el camino seguido por los paquetes.

\begin{longtable}{|p{2.8cm}|p{3cm}|p{4cm}|p{5.6cm}|p{1.5cm}|}

\hline
Comando &
Sintaxis &
Argumentos útiles &
Particularidad &
OS
\\
\hline
\endhead

traceroute &
traceroute destino &
-I, -T, -m, -w &
Obtiene cada uno de los saltos recorridos hasta el destino utilizando paquetes ICMP, UDP o TCP. Fundamental para detectar dónde se pierde una conexión. &
Linux/macOS
\\
\hline

tracert &
tracert destino &
-d, -h, -w &
Equivalente disponible en Windows. El argumento \texttt{-d} evita resolver nombres DNS y acelera considerablemente la ejecución. &
Windows
\\
\hline

tracepath &
tracepath destino &
N/A &
Similar a traceroute pero sin requerir privilegios administrativos. Además detecta automáticamente el MTU del camino. &
Linux
\\
\hline

mtr &
mtr destino &
-r, -c, -w &
Combina las funcionalidades de ping y traceroute mostrando estadísticas en tiempo real sobre cada salto. Una de las herramientas más utilizadas para diagnosticar enlaces inestables. &
Linux/macOS
\\
\hline

pathping &
pathping destino &
-n &
Herramienta de Windows que combina traceroute y estadísticas de pérdida de paquetes por salto. Muy útil para identificar enlaces congestionados. &
Windows
\\
\hline

\end{longtable}


%%%%%%%%%%%%%%%55

%%%%%%%%%%%%%%%%%%%%%%%%%%%%%%%%%%%%%%%%%%%%%%%%%%%%%%%%%%%%%%%%%%%%%%%%%%%%%%%

\subsubsection{Verificación de puertos remotos}

En numerosas ocasiones un equipo responde correctamente a ICMP, pero el puerto
utilizado por la aplicación permanece cerrado o inaccesible.

Los siguientes comandos permiten comprobar si un puerto TCP o UDP acepta
conexiones.

\begin{longtable}{|p{2.8cm}|p{3cm}|p{4cm}|p{5.6cm}|p{1.5cm}|}

\hline
Comando &
Sintaxis &
Argumentos útiles &
Particularidad &
OS
\\
\hline
\endhead

nc &
nc host puerto &
-z, -v, -u, -w &
Netcat permite comprobar rápidamente si un puerto remoto acepta conexiones TCP o UDP. Constituye una de las herramientas más utilizadas durante troubleshooting de aplicaciones. &
Linux/macOS
\\
\hline

netcat &
netcat host puerto &
-zv &
Nombre completo de la utilidad en algunas distribuciones. &
Linux
\\
\hline

telnet &
telnet host puerto &
N/A &
Permite abrir manualmente una conexión TCP hacia un puerto determinado. Aunque es una herramienta antigua, continúa siendo útil para comprobar disponibilidad de servicios. &
Todos
\\
\hline

Test-NetConnection &
Test-NetConnection host -Port puerto &
-InformationLevel &
Herramienta de PowerShell que verifica conectividad TCP, resolución DNS y rutas de manera integrada. Equivale funcionalmente a combinar ping y telnet. &
Windows
\\
\hline

curl &
curl host:puerto &
-v &
Aunque está orientado a protocolos de aplicación, también permite comprobar rápidamente si un servicio HTTP o HTTPS responde. Se estudiará con mayor detalle en la sección correspondiente. &
Todos
\\
\hline

wget &
wget URL &
-S, --spider &
Permite verificar si un recurso HTTP se encuentra disponible sin descargar completamente el contenido. &
Linux
\\
\hline

\end{longtable}


%%%%%%%%%%%%%%5
%%%%%%%%%%%%%%%%%%%%%%%%%%%%%%%%%%%%%%%%%%%%%%%%%%%%%%%%%%%%%%%%%%%%%%%%%%%%%%%

\subsubsection{Disponibilidad de hosts}

En algunas situaciones resulta necesario determinar rápidamente qué equipos de
una red se encuentran activos.

\begin{longtable}{|p{2.8cm}|p{3cm}|p{4cm}|p{5.6cm}|p{1.5cm}|}

\hline
Comando &
Sintaxis &
Argumentos útiles &
Particularidad &
OS
\\
\hline
\endhead

ping &
for ip in ...; do ping ... &
-c 1 &
Puede utilizarse dentro de scripts para comprobar múltiples equipos. &
Linux/macOS
\\
\hline

fping &
fping -a -g red/24 &
-g &
Descubre rápidamente qué direcciones IP responden dentro de un segmento de red. &
Linux
\\
\hline

arp-scan &
arp-scan red &
--localnet &
Realiza un descubrimiento de hosts mediante ARP dentro de la red local. Funciona incluso cuando ICMP está filtrado. &
Linux
\\
\hline

nmap &
nmap -sn red &
-PR, -PE &
Realiza descubrimiento de hosts sin escanear puertos. Más adelante se profundizará en esta herramienta. &
Linux/macOS/Windows
\\
\hline

\end{longtable}


%%%%%%

%%%%%%%%%%%%%%%%%%%%%%%%%%%%%%%%%%%%%%%%%%%%%%%%%%%%%%%%%%%%%%%%%%%%%%%%%%%%%%%

\subsection{Paso 4. Verificación de resolución DNS}

Una vez comprobada la conectividad IP, el siguiente paso consiste en verificar
la correcta resolución de nombres.

En muchas ocasiones un servicio se encuentra disponible y accesible mediante su
dirección IP, pero falla al utilizar su nombre de dominio debido a problemas de
configuración DNS.

Durante esta etapa se busca responder preguntas como:

\begin{itemize}
\item ¿El nombre realmente resuelve?
\item ¿Qué servidor DNS respondió?
\item ¿Cuál fue la dirección IP obtenida?
\item ¿La respuesta proviene de caché?
\item ¿Los registros DNS son los esperados?
\item ¿Existe algún problema de propagación?
\end{itemize}

%%%%%%%%%%%%%%%%%%%%%%%%%%%%%%%%%%%%%%%%%%%%%%%%%%%%%%%%%%%%%%%%%%%%%%%%%%%%%%%

\subsubsection{Consultas DNS}

Las siguientes herramientas permiten consultar registros DNS directamente desde
uno o varios servidores.

\begin{longtable}{|p{2.8cm}|p{3cm}|p{4cm}|p{5.6cm}|p{1.5cm}|}

\hline
Comando &
Sintaxis &
Argumentos útiles &
Particularidad &
OS
\\
\hline
\endhead

nslookup &
nslookup dominio.com &
server, set type, exit &
Permite realizar consultas DNS especificando el servidor de nombres y el tipo de registro a consultar. Disponible prácticamente en todos los sistemas operativos. &
Todos
\\
\hline

dig &
dig dominio.com &
+A, +short, +trace, +tcp &
Herramienta estándar para consultas DNS avanzadas. Permite inspeccionar respuestas completas, seguir la delegación desde los servidores raíz y realizar consultas TCP cuando sea necesario. &
Linux/macOS
\\
\hline

host &
host dominio.com &
-t, -a &
Realiza consultas DNS rápidas para distintos tipos de registros. Muy útil para verificaciones sencillas. &
Linux/macOS
\\
\hline

Resolve-DnsName &
Resolve-DnsName dominio.com &
-Type, -Server &
Herramienta de PowerShell que permite consultar registros DNS especificando el tipo de registro y el servidor DNS a utilizar. &
Windows
\\
\hline

getent hosts &
getent hosts dominio.com &
ahosts &
Consulta la resolución de nombres utilizando los mecanismos configurados por el sistema operativo (DNS, archivos hosts, LDAP, etc.). Muy útil para verificar el comportamiento real del resolver del sistema. &
Linux
\\
\hline

dscacheutil &
dscacheutil -q host -a name dominio.com &
N/A &
Consulta la resolución de nombres utilizando la caché del sistema en macOS. &
macOS
\\
\hline

\end{longtable}

%%%%%%%%%%%%%%%%%%%%%%%%%%%%%%%%%%%%%%%%%%%%%%%%%%%%%%%%%%%%%%%%%%%%%%%%%%%%%%%

\subsubsection{Tipos de registros DNS}

Durante el troubleshooting suele ser necesario consultar registros específicos.

\begin{longtable}{|p{2.8cm}|p{3cm}|p{4cm}|p{5.6cm}|p{1.5cm}|}

\hline
Registro &
Consulta &
Uso principal &
Descripción &
OS
\\
\hline
\endhead

A &
dig dominio A &
IPv4 &
Obtiene la dirección IPv4 asociada a un nombre de dominio. &
Todos
\\
\hline

AAAA &
dig dominio AAAA &
IPv6 &
Obtiene la dirección IPv6 correspondiente. &
Todos
\\
\hline

MX &
dig dominio MX &
Correo electrónico &
Consulta los servidores responsables de recibir correo electrónico para un dominio. &
Todos
\\
\hline

NS &
dig dominio NS &
Delegación &
Obtiene los servidores autoritativos del dominio. &
Todos
\\
\hline

SOA &
dig dominio SOA &
Administración &
Consulta el registro Start Of Authority del dominio. Muy útil para revisar números de serie y temporizadores de zona. &
Todos
\\
\hline

TXT &
dig dominio TXT &
Validaciones &
Permite consultar registros TXT utilizados para SPF, DKIM, DMARC y diversas verificaciones. &
Todos
\\
\hline

CNAME &
dig dominio CNAME &
Alias &
Obtiene alias canónicos de un nombre de dominio. &
Todos
\\
\hline

PTR &
dig -x direcciónIP &
DNS inverso &
Permite realizar consultas de resolución inversa. Muy utilizado durante troubleshooting de correo electrónico. &
Todos
\\
\hline

SRV &
dig dominio SRV &
Servicios &
Consulta registros utilizados por servicios como Active Directory, SIP o Kerberos. &
Todos
\\
\hline

\end{longtable}

%%%%%%%%%%%%%%%%%%%%%%%%%%%%%%%%%%%%%%%%%%%%%%%%%%%%%%%%%%%%%%%%%%%%%%%%%%%%%%%

\subsubsection{Diagnóstico del resolver local}

En algunas ocasiones el problema no reside en el servidor DNS sino en la
configuración local del resolver.

Las siguientes herramientas permiten inspeccionar dicha configuración.

\begin{longtable}{|p{2.8cm}|p{3cm}|p{4cm}|p{5.6cm}|p{1.5cm}|}

\hline
Comando &
Sintaxis &
Argumentos útiles &
Particularidad &
OS
\\
\hline
\endhead

cat &
cat /etc/resolv.conf &
N/A &
Permite consultar los servidores DNS configurados por el sistema. &
Linux
\\
\hline

resolvectl &
resolvectl status &
query &
Consulta el estado del resolver administrado por systemd-resolved y los servidores DNS asociados a cada interfaz. &
Linux
\\
\hline

systemd-resolve &
systemd-resolve --status &
--statistics &
Versión utilizada en algunas distribuciones para consultar el estado del resolver. &
Linux
\\
\hline

nmcli &
nmcli device show &
connection show &
Permite consultar los servidores DNS configurados por NetworkManager. &
Linux
\\
\hline

scutil &
scutil --dns &
N/A &
Muestra la configuración completa del resolver DNS en macOS. &
macOS
\\
\hline

ipconfig &
ipconfig /all &
N/A &
Permite consultar los servidores DNS configurados en Windows. &
Windows
\\
\hline

Get-DnsClientServerAddress &
Get-DnsClientServerAddress &
InterfaceAlias &
Obtiene la configuración DNS mediante PowerShell. &
Windows
\\
\hline

\end{longtable}

%%%%%%%%%%%%%%%%%%%%%%%%%%%%%%%%%%%%%%%%%%%%%%%%%%%%%%%%%%%%%%%%%%%%%%%%%%%%%%%

\subsubsection{Caché DNS}

En determinadas ocasiones el sistema operativo mantiene información DNS
almacenada localmente. Si un registro cambió recientemente o la información
almacenada es incorrecta, el equipo puede continuar utilizando una dirección
IP obsoleta.

Las siguientes herramientas permiten inspeccionar o limpiar la caché DNS.

\begin{longtable}{|p{2.8cm}|p{3cm}|p{4cm}|p{5.6cm}|p{1.5cm}|}

\hline
Comando &
Sintaxis &
Argumentos útiles &
Particularidad &
OS
\\
\hline
\endhead

ipconfig &
ipconfig /displaydns &
N/A &
Muestra el contenido de la caché DNS local de Windows. Muy útil para verificar si un registro se encuentra almacenado localmente. &
Windows
\\
\hline

ipconfig &
ipconfig /flushdns &
N/A &
Vacía completamente la caché DNS local. Es uno de los comandos más utilizados durante troubleshooting. &
Windows
\\
\hline

Resolve-DnsName &
Resolve-DnsName dominio &
-CacheOnly &
Consulta únicamente la información almacenada en la caché DNS local. &
Windows
\\
\hline

resolvectl &
resolvectl statistics &
reset-statistics &
Permite consultar estadísticas del resolver administrado por systemd. &
Linux
\\
\hline

resolvectl &
resolvectl flush-caches &
N/A &
Vacía la caché administrada por systemd-resolved. &
Linux
\\
\hline

systemd-resolve &
systemd-resolve --statistics &
--flush-caches &
Versión disponible en algunas distribuciones anteriores a resolvectl. &
Linux
\\
\hline

dscacheutil &
dscacheutil -flushcache &
N/A &
Vacía la caché DNS del sistema en macOS. &
macOS
\\
\hline

killall &
killall -HUP mDNSResponder &
N/A &
Solicita al servicio mDNSResponder reconstruir la caché DNS. Método habitual para limpiar la caché DNS en macOS. &
macOS
\\
\hline

\end{longtable}

%%%%%%%%%%%%%%%%%%%%%%%%%%%%%%%%%%%%%%%%%%%%%%%%%%%%%%%%%%%%%%%%%%%%%%%%%%%%%%%

\subsubsection{Renovación de configuración DHCP}

Cuando un equipo obtiene su dirección IP mediante DHCP, una concesión vencida
o una configuración incorrecta pueden impedir la comunicación con otros equipos.

Los siguientes comandos permiten renovar dicha configuración.

\begin{longtable}{|p{2.8cm}|p{3cm}|p{4cm}|p{5.6cm}|p{1.5cm}|}

\hline
Comando &
Sintaxis &
Argumentos útiles &
Particularidad &
OS
\\
\hline
\endhead

ipconfig &
ipconfig /release &
N/A &
Libera la dirección IP obtenida mediante DHCP. &
Windows
\\
\hline

ipconfig &
ipconfig /renew &
N/A &
Solicita una nueva concesión DHCP al servidor correspondiente. &
Windows
\\
\hline

dhclient &
dhclient interfaz &
-r &
Cliente DHCP utilizado por numerosas distribuciones Linux para renovar direcciones IP. &
Linux
\\
\hline

nmcli &
nmcli connection up conexión &
down &
Permite reiniciar una conexión administrada por NetworkManager, solicitando nuevamente la configuración DHCP cuando corresponda. &
Linux
\\
\hline

networksetup &
networksetup -renewdhcp interfaz &
N/A &
Renueva la concesión DHCP en macOS. &
macOS
\\
\hline

\end{longtable}

%%%%%%%%%%%%%%%%%%%%%%%%%%%%%%%%%%%%%%%%%%%%%%%%%%%%%%%%%%%%%%%%%%%%%%%%%%%%%%%

\subsubsection{Archivo hosts}

Antes de consultar un servidor DNS, la mayoría de los sistemas operativos
verifican si el nombre solicitado se encuentra definido localmente en el
archivo hosts.

Una entrada incorrecta puede provocar que un dominio resuelva hacia una
dirección IP equivocada.

\begin{longtable}{|p{2.8cm}|p{3cm}|p{4cm}|p{5.6cm}|p{1.5cm}|}

\hline
Comando &
Sintaxis &
Argumentos útiles &
Particularidad &
OS
\\
\hline
\endhead

cat &
cat /etc/hosts &
N/A &
Permite consultar las entradas definidas en el archivo hosts. &
Linux/macOS
\\
\hline

less &
less /etc/hosts &
N/A &
Facilita revisar archivos hosts extensos. &
Linux/macOS
\\
\hline

grep &
grep nombre /etc/hosts &
-i &
Busca rápidamente una entrada específica dentro del archivo hosts. &
Linux/macOS
\\
\hline

type &
type C:\\Windows\\System32\\drivers\\etc\\hosts &
N/A &
Muestra el contenido del archivo hosts en Windows. &
Windows
\\
\hline

findstr &
findstr dominio hosts &
/I &
Busca una entrada específica dentro del archivo hosts. &
Windows
\\
\hline

\end{longtable}


%%%%%%%%%%%

%%%%%%%%%%%%%%%%%%%%%%%%%%%%%%%%%%%%%%%%%%%%%%%%%%%%%%%%%%%%%%%%%%%%%%%%%%%%%%%

\subsection{Paso 5. Verificación de puertos y conexiones}

Una vez comprobada la resolución DNS, el siguiente paso consiste en verificar
si el servicio realmente se encuentra escuchando conexiones y si el sistema
permite establecer comunicación con él.

En esta etapa se busca responder preguntas como:

\begin{itemize}
\item ¿El proceso realmente abrió el puerto esperado?
\item ¿Qué proceso está utilizando dicho puerto?
\item ¿Existe alguna conexión establecida?
\item ¿El puerto está en estado LISTEN?
\item ¿Existe algún conflicto de puertos?
\end{itemize}

%%%%%%%%%%%%%%%%%%%%%%%%%%%%%%%%%%%%%%%%%%%%%%%%%%%%%%%%%%%%%%%%%%%%%%%%%%%%%%%

\subsubsection{Puertos en escucha y conexiones activas}

Las siguientes herramientas permiten inspeccionar conexiones TCP, UDP y
puertos abiertos tanto locales como remotos.

\begin{longtable}{|p{2.8cm}|p{3cm}|p{4cm}|p{5.6cm}|p{1.5cm}|}

\hline
Comando &
Sintaxis &
Argumentos útiles &
Particularidad &
OS
\\
\hline
\endhead

ss &
ss -tulpn &
-a, -l, -n, -p &
Herramienta moderna para inspeccionar conexiones TCP/UDP y puertos en escucha. Sustituye progresivamente a netstat en Linux. &
Linux
\\
\hline

ss &
ss -tan &
state established &
Permite visualizar conexiones TCP activas y filtrar por estado. &
Linux
\\
\hline

netstat &
netstat -tulpn &
-r, -a, -n, -p &
Herramienta clásica para visualizar conexiones, interfaces y puertos abiertos. Continúa disponible en numerosos sistemas. &
Linux/macOS
\\
\hline

netstat &
netstat -ano &
-b &
Permite visualizar conexiones, PID y puertos abiertos en Windows. &
Windows
\\
\hline

Get-NetTCPConnection &
Get-NetTCPConnection &
-State &
Obtiene todas las conexiones TCP activas mediante PowerShell. &
Windows
\\
\hline

Get-NetUDPEndpoint &
Get-NetUDPEndpoint &
N/A &
Lista los endpoints UDP abiertos. &
Windows
\\
\hline

lsof &
lsof -i &
-P, -n &
Relaciona conexiones de red con los procesos que las mantienen abiertas. Constituye una de las herramientas más útiles durante troubleshooting. &
Linux/macOS
\\
\hline

lsof &
lsof -i :443 &
:80 &
Permite conocer qué proceso está utilizando un puerto específico. &
Linux/macOS
\\
\hline

fuser &
fuser 443/tcp &
-v, -k &
Identifica el proceso asociado a un puerto y puede finalizarlo si es necesario. &
Linux
\\
\hline


%%%%
\hline

sockstat &
sockstat &
-4, -6, -l &
Permite visualizar sockets abiertos y los procesos asociados. Es el equivalente de \texttt{lsof} en sistemas BSD y macOS. &
macOS
\\
\hline

netstat &
netstat -s &
-t, -u &
Presenta estadísticas detalladas de los protocolos TCP, UDP, ICMP e IP. Muy útil para detectar retransmisiones, errores de checksum o paquetes descartados. &
Linux/macOS
\\
\hline

Get-NetTCPConnection &
Get-NetTCPConnection -State Listen &
LocalPort &
Permite listar únicamente los puertos que se encuentran escuchando conexiones. &
Windows
\\
\hline

Get-NetUDPEndpoint &
Get-NetUDPEndpoint &
LocalPort &
Consulta todos los sockets UDP abiertos mediante PowerShell. &
Windows
\\
\hline

\end{longtable}

%%%%%%%%%%%%%%%%%%%%%%%%%%%%%%%%%%%%%%%%%%%%%%%%%%%%%%%%%%%%%%%%%%%%%%%%%%%%%%%

\subsubsection{Procesos asociados a conexiones}

Una vez identificado un puerto o una conexión, suele ser necesario determinar
qué proceso la abrió.

Las siguientes herramientas permiten relacionar conexiones de red con procesos.

\begin{longtable}{|p{2.8cm}|p{3cm}|p{4cm}|p{5.6cm}|p{1.5cm}|}

\hline
Comando &
Sintaxis &
Argumentos útiles &
Particularidad &
OS
\\
\hline
\endhead

lsof &
lsof -iTCP &
-sTCP:LISTEN &
Permite listar únicamente procesos que mantienen sockets TCP en estado LISTEN. Muy útil para confirmar que una aplicación realmente abrió el puerto esperado. &
Linux/macOS
\\
\hline

lsof &
lsof -iTCP -sTCP:ESTABLISHED &
-nP &
Lista conexiones TCP actualmente establecidas junto con el proceso correspondiente. &
Linux/macOS
\\
\hline

lsof &
lsof -p PID &
N/A &
Permite conocer todos los archivos, sockets y recursos abiertos por un proceso determinado. &
Linux/macOS
\\
\hline

fuser &
fuser puerto/tcp &
-v &
Obtiene el proceso que utiliza un puerto específico. &
Linux
\\
\hline

fuser &
fuser archivo &
-v &
Permite determinar qué proceso mantiene abierto un archivo. Muy útil cuando un archivo no puede eliminarse. &
Linux
\\
\hline

netstat &
netstat -ano &
N/A &
Relaciona conexiones TCP con el PID correspondiente. &
Windows
\\
\hline

tasklist &
tasklist /FI "PID eq xxxx" &
/SVC &
Permite identificar el proceso correspondiente a un PID obtenido mediante netstat. &
Windows
\\
\hline

Get-Process &
Get-Process -Id PID &
N/A &
Obtiene información adicional del proceso asociado a una conexión. &
Windows
\\
\hline

\end{longtable}

%%%%%%%%%%%%%%%%%%%%%%%%%%%%%%%%%%%%%%%%%%%%%%%%%%%%%%%%%%%%%%%%%%%%%%%%%%%%%%%

\subsection{Paso 6. Verificación del servicio HTTP}

Una vez comprobado que existe conectividad TCP hacia el puerto correspondiente,
el siguiente paso consiste en verificar que la aplicación realmente responda
utilizando el protocolo esperado.

Durante esta etapa se busca responder preguntas como:

\begin{itemize}
\item ¿El servidor responde?
\item ¿Qué código HTTP devuelve?
\item ¿Existen redirecciones?
\item ¿Qué encabezados devuelve?
\item ¿Qué versión HTTP utiliza?
\item ¿Cuál fue el tiempo de respuesta?
\end{itemize}

%%%%%%%%%%%%%%%%%%%%%%%%%%%%%%%%%%%%%%%%%%%%%%%%%%%%%%%%%%%%%%%%%%%%%%%%%%%%%%%

\subsubsection{Consultas HTTP}

Las siguientes herramientas permiten enviar solicitudes HTTP y analizar la
respuesta del servidor.

\begin{longtable}{|p{2.8cm}|p{3cm}|p{4cm}|p{5.6cm}|p{1.5cm}|}

\hline
Comando &
Sintaxis &
Argumentos útiles &
Particularidad &
OS
\\
\hline
\endhead

curl &
curl URL &
-v, -I, -L, -O, -o &
Herramienta estándar para realizar solicitudes HTTP y HTTPS. Permite inspeccionar encabezados, seguir redirecciones y descargar recursos. Constituye una de las utilidades más importantes durante troubleshooting. &
Todos
\\
\hline

curl &
curl -I URL &
N/A &
Solicita únicamente los encabezados HTTP utilizando el método HEAD. Muy útil para comprobar rápidamente la disponibilidad de un servicio. &
Todos
\\
\hline

curl &
curl -v URL &
--http1.1, --http2 &
Muestra información detallada de la conexión, incluyendo encabezados enviados y recibidos. &
Todos
\\
\hline

curl &
curl -L URL &
--max-redirs &
Sigue automáticamente redirecciones HTTP. &
Todos
\\
\hline

curl &
curl -X POST URL &
-H, -d &
Permite probar distintos métodos HTTP enviando encabezados y cuerpos personalizados. &
Todos
\\
\hline

curl &
curl --resolve dominio:443:IP &
--connect-to &
Permite forzar la resolución DNS hacia una dirección IP específica sin modificar el archivo hosts. Muy útil durante migraciones y pruebas. &
Todos
\\
\hline

wget &
wget URL &
-S, --spider, -O &
Permite descargar recursos o verificar su existencia sin almacenarlos localmente. &
Linux
\\
\hline

wget &
wget --spider URL &
-S &
Comprueba la disponibilidad de un recurso HTTP simulando una descarga. &
Linux
\\
\hline

Invoke-WebRequest &
Invoke-WebRequest URL &
-Method, -Headers &
Herramienta de PowerShell para realizar solicitudes HTTP y HTTPS. &
Windows
\\
\hline

Invoke-RestMethod &
Invoke-RestMethod URL &
-Method &
Orientada principalmente a servicios REST y APIs. Convierte automáticamente respuestas JSON y XML en objetos de PowerShell. &
Windows
\\
\hline

\end{longtable}

%%%%%%
%%%%%%%%%%%%%%%%%%%%%%%%%%%%%%%%%%%%%%%%%%%%%%%%%%%%%%%%%%%%%%%%%%%%%%%%%%%%%%%

\subsection{Paso 7. Verificación de HTTPS y TLS}

Una vez validado que el servidor responde mediante HTTP o HTTPS, resulta
necesario verificar el establecimiento de la conexión TLS.

Muchos problemas de conectividad corresponden realmente a errores durante el
handshake TLS, certificados expirados, cadenas de confianza incompletas,
versiones incompatibles del protocolo o configuraciones incorrectas del
servidor.

Durante esta etapa se busca responder preguntas como:

\begin{itemize}
\item ¿El handshake TLS se completa correctamente?
\item ¿El certificado es válido?
\item ¿La cadena de certificados está completa?
\item ¿Qué versión de TLS fue negociada?
\item ¿Qué Cipher Suite se seleccionó?
\item ¿El servidor requiere SNI?
\item ¿El certificado corresponde al dominio solicitado?
\end{itemize}

%%%%%%%%%%%%%%%%%%%%%%%%%%%%%%%%%%%%%%%%%%%%%%%%%%%%%%%%%%%%%%%%%%%%%%%%%%%%%%%

\subsubsection{Inspección del Handshake TLS}

Las siguientes herramientas permiten inspeccionar el establecimiento de una
conexión TLS y obtener información detallada sobre el certificado presentado
por el servidor.

\begin{longtable}{|p{2.8cm}|p{3cm}|p{4cm}|p{5.6cm}|p{1.5cm}|}

\hline
Comando &
Sintaxis &
Argumentos útiles &
Particularidad &
OS
\\
\hline
\endhead

openssl s\_client &
openssl s\_client -connect dominio:443 &
-showcerts, -servername, -brief &
Herramienta principal para diagnosticar conexiones TLS. Permite visualizar el handshake completo, la cadena de certificados y los parámetros negociados. &
Linux/macOS
\\
\hline

openssl s\_client &
openssl s\_client -connect dominio:443 -servername dominio &
-showcerts &
Permite enviar el nombre del servidor mediante SNI. Fundamental cuando un mismo servidor hospeda múltiples sitios HTTPS. &
Linux/macOS
\\
\hline

openssl s\_client &
openssl s\_client -connect dominio:443 -tls1\_2 &
-tls1\_3 &
Permite forzar una versión específica del protocolo TLS durante la negociación. &
Linux/macOS
\\
\hline

openssl s\_client &
openssl s\_client -connect dominio:443 -cipher 'ECDHE:*' &
-cipher &
Permite comprobar qué Cipher Suites acepta el servidor. &
Linux/macOS
\\
\hline

curl &
curl -v https://dominio &
--http1.1, --http2 &
Muestra información detallada de la conexión HTTPS, incluyendo negociación TLS y certificados utilizados. &
Todos
\\
\hline

curl &
curl --tlsv1.2 URL &
--tlsv1.3 &
Fuerza el uso de una versión determinada de TLS para comprobar compatibilidad. &
Todos
\\
\hline

\end{longtable}

%%%%%%%%%%%%%%%%%%%%%%%%%%%%%%%%%%%%%%%%%%%%%%%%%%%%%%%%%%%%%%%%%%%%%%%%%%%%%%%

\subsubsection{Inspección de certificados}

Una vez obtenido el certificado del servidor, resulta necesario verificar
su validez, fechas, emisor y sujeto.

\begin{longtable}{|p{2.8cm}|p{3cm}|p{4cm}|p{5.6cm}|p{1.5cm}|}

\hline
Comando &
Sintaxis &
Argumentos útiles &
Particularidad &
OS
\\
\hline
\endhead

openssl x509 &
openssl x509 -in certificado.pem -text &
-noout &
Muestra toda la información contenida en un certificado X.509. &
Linux/macOS
\\
\hline

openssl x509 &
openssl x509 -dates &
-noout &
Obtiene las fechas de inicio y expiración del certificado. Muy útil para detectar certificados vencidos. &
Linux/macOS
\\
\hline

openssl x509 &
openssl x509 -subject &
-issuer &
Permite consultar el sujeto y la autoridad certificadora que emitió el certificado. &
Linux/macOS
\\
\hline

openssl x509 &
openssl x509 -fingerprint &
-sha256 &
Obtiene la huella digital del certificado para comparaciones de integridad. &
Linux/macOS
\\
\hline

openssl verify &
openssl verify certificado.pem &
-CAfile &
Comprueba si un certificado puede validarse correctamente contra una autoridad certificadora determinada. &
Linux/macOS
\\
\hline

certutil &
certutil -verify certificado.cer &
-dump &
Permite validar certificados y visualizar su contenido desde Windows. &
Windows
\\
\hline

\end{longtable}


%%%%%%%%%%%%%%%%%%%%%%%%%%%%%%%%%%%%%%%%%%%%%%%%%%%%%%%%%%%%%%%%%%%%%%%%%%%%%%%

\subsubsection{Pruebas de autenticación TLS}

En algunos servicios HTTPS el cliente también debe presentar un certificado
durante el establecimiento de la conexión.

Los siguientes comandos permiten realizar este tipo de pruebas.

\begin{longtable}{|p{2.8cm}|p{3cm}|p{4cm}|p{5.6cm}|p{1.5cm}|}

\hline
Comando &
Sintaxis &
Argumentos útiles &
Particularidad &
OS
\\
\hline
\endhead

curl &
curl --cert cliente.pem URL &
--key &
Permite autenticarse utilizando certificados de cliente. Muy utilizado en APIs empresariales y autenticación mutua (mTLS). &
Todos
\\
\hline

curl &
curl --cacert ca.pem URL &
--capath &
Especifica explícitamente la autoridad certificadora utilizada para validar el servidor. &
Todos
\\
\hline

curl &
curl --insecure URL &
-k &
Omite la validación del certificado del servidor. Debe utilizarse únicamente para fines de diagnóstico. &
Todos
\\
\hline

openssl s\_client &
openssl s\_client -cert cliente.pem -key cliente.key &
-CAfile &
Permite realizar pruebas completas de autenticación mutua utilizando OpenSSL. &
Linux/macOS
\\
\hline

\end{longtable}


%%%%%%%%%%%%%%%%%%%%%%%%%%%%%%%%%%%%%%%%%%%%%%%%%%%%%%%%%%%%%%%%%%%%%%%%%%%%%%%

\subsubsection{Información sobre Cipher Suites y protocolos}

Durante el troubleshooting también resulta útil conocer qué versiones de TLS
y qué algoritmos criptográficos soporta una instalación determinada.

\begin{longtable}{|p{2.8cm}|p{3cm}|p{4cm}|p{5.6cm}|p{1.5cm}|}

\hline
Comando &
Sintaxis &
Argumentos útiles &
Particularidad &
OS
\\
\hline
\endhead

openssl version &
openssl version &
-a &
Muestra la versión instalada de OpenSSL y las características compiladas. &
Linux/macOS
\\
\hline

openssl ciphers &
openssl ciphers &
-v, -tls1\_2, -tls1\_3 &
Lista las Cipher Suites soportadas por la instalación local de OpenSSL. &
Linux/macOS
\\
\hline

openssl list &
openssl list -cipher-algorithms &
-digest-algorithms &
Permite consultar los algoritmos criptográficos disponibles. &
Linux/macOS
\\
\hline

curl &
curl --http2 URL &
--http1.1 &
Permite comprobar la versión HTTP negociada con el servidor. &
Todos
\\
\hline

curl &
curl --tls-max 1.2 URL &
--tlsv1.3 &
Restringe la versión máxima del protocolo TLS utilizada durante la conexión. Muy útil para detectar incompatibilidades de versión. &
Todos
\\
\hline

\end{longtable}

%%%%%%%%%%%%%%%%%%%%%%%%%%%%%%%%%%%%%%%%%%%%%%%%%%%%%%%%%%%%%%%%%%%%%%%%%%%%%%%

\subsection{Paso 8. Verificación de Firewalls}

Una vez comprobada la conectividad, la resolución DNS y la disponibilidad del
servicio, resulta necesario verificar que ningún firewall esté bloqueando la
comunicación.

Durante esta etapa se busca responder preguntas como:

\begin{itemize}
\item ¿Existe un firewall habilitado?
\item ¿Qué política predeterminada utiliza?
\item ¿Existe una regla que bloquee el puerto?
\item ¿El tráfico está siendo descartado (DROP) o rechazado (REJECT)?
\item ¿La regla realmente está siendo utilizada?
\end{itemize}

%%%%%%%%%%%%%%%%%%%%%%%%%%%%%%%%%%%%%%%%%%%%%%%%%%%%%%%%%%%%%%%%%%%%%%%%%%%%%%%

\subsubsection{iptables}

Aunque numerosas distribuciones modernas utilizan \texttt{nftables}, todavía
es muy frecuente encontrar sistemas administrados mediante \texttt{iptables}.

\begin{longtable}{|p{2.8cm}|p{3cm}|p{4cm}|p{5.6cm}|p{1.5cm}|}

\hline
Comando &
Sintaxis &
Argumentos útiles &
Particularidad &
OS
\\
\hline
\endhead

iptables &
iptables -L &
-v, -n, --line-numbers &
Lista las reglas configuradas mostrando contadores de paquetes y bytes. Constituye normalmente el primer comando durante troubleshooting. &
Linux
\\
\hline

iptables &
iptables -S &
N/A &
Muestra las reglas utilizando la sintaxis completa de iptables. Muy útil para reconstruir configuraciones. &
Linux
\\
\hline

iptables &
iptables -t nat -L &
-v &
Consulta específicamente las reglas de NAT. &
Linux
\\
\hline

iptables &
iptables -t mangle -L &
-v &
Consulta reglas de modificación de paquetes. &
Linux
\\
\hline

iptables-save &
iptables-save &
N/A &
Exporta completamente la configuración actual del firewall. Muy útil antes de realizar modificaciones. &
Linux
\\
\hline

iptables &
iptables -F &
N/A &
Elimina todas las reglas de una cadena. Debe utilizarse con extrema precaución. &
Linux
\\
\hline

iptables &
iptables -A INPUT ... &
-A, -I, -D &
Permite agregar, insertar o eliminar reglas específicas. &
Linux
\\
\hline

iptables &
iptables -P INPUT ACCEPT &
DROP, FORWARD, OUTPUT &
Consulta o modifica la política predeterminada de una cadena. &
Linux
\\
\hline

\end{longtable}

%%%%%%%%%%%%%%%%%%%%%%%%%%%%%%%%%%%%%%%%%%%%%%%%%%%%%%%%%%%%%%%%%%%%%%%%%%%%%%%

\subsubsection{nftables}

Las distribuciones Linux actuales utilizan cada vez con mayor frecuencia
\texttt{nftables} como reemplazo de \texttt{iptables}.

\begin{longtable}{|p{2.8cm}|p{3cm}|p{4cm}|p{5.6cm}|p{1.5cm}|}

\hline
Comando &
Sintaxis &
Argumentos útiles &
Particularidad &
OS
\\
\hline
\endhead

nft &
nft list ruleset &
N/A &
Muestra todas las tablas, cadenas y reglas configuradas en nftables. &
Linux
\\
\hline

nft &
nft list table inet filter &
N/A &
Consulta únicamente una tabla específica. &
Linux
\\
\hline

nft &
nft list chain inet filter INPUT &
N/A &
Visualiza las reglas pertenecientes a una cadena determinada. &
Linux
\\
\hline

nft &
nft add rule ... &
delete rule &
Permite agregar o eliminar reglas del firewall. &
Linux
\\
\hline

nft &
nft monitor &
trace &
Monitorea modificaciones realizadas sobre las reglas del firewall en tiempo real. &
Linux
\\
\hline

\end{longtable}

%%%%%%%%%%%%%%%%%%%%%%%%%%%%%%%%%%%%%%%%%%%%%%%%%%%%%%%%%%%%%%%%%%%%%%%%%%%%%%%

\subsubsection{UFW}

Ubuntu y otras distribuciones utilizan con frecuencia UFW como interfaz
simplificada sobre iptables.

\begin{longtable}{|p{2.8cm}|p{3cm}|p{4cm}|p{5.6cm}|p{1.5cm}|}

\hline
Comando &
Sintaxis &
Argumentos útiles &
Particularidad &
OS
\\
\hline
\endhead

ufw &
ufw status &
verbose, numbered &
Consulta el estado actual del firewall y las reglas configuradas. &
Linux
\\
\hline

ufw &
ufw allow puerto &
deny, reject &
Permite habilitar o bloquear puertos específicos. &
Linux
\\
\hline

ufw &
ufw delete número &
allow, deny &
Elimina reglas previamente configuradas. &
Linux
\\
\hline

ufw &
ufw reload &
N/A &
Recarga la configuración del firewall. &
Linux
\\
\hline

ufw &
ufw enable &
disable &
Habilita o deshabilita completamente UFW. &
Linux
\\
\hline

\end{longtable}


%%%%%%


%%%%%%%%%%%%%%%%%%%%%%%%%%%%%%%%%%%%%%%%%%%%%%%%%%%%%%%%%%%%%%%%%%%%%%%%%%%%%%%

\subsubsection{firewalld}

Algunas distribuciones empresariales como Red Hat Enterprise Linux,
CentOS, Rocky Linux o AlmaLinux utilizan \texttt{firewalld}.

\begin{longtable}{|p{2.8cm}|p{3cm}|p{4cm}|p{5.6cm}|p{1.5cm}|}

\hline
Comando &
Sintaxis &
Argumentos útiles &
Particularidad &
OS
\\
\hline
\endhead

firewall-cmd &
firewall-cmd --state &
N/A &
Consulta si firewalld se encuentra activo. &
Linux
\\
\hline

firewall-cmd &
firewall-cmd --list-all &
--zone &
Muestra la configuración completa de una zona. &
Linux
\\
\hline

firewall-cmd &
firewall-cmd --list-ports &
N/A &
Lista todos los puertos abiertos. &
Linux
\\
\hline

firewall-cmd &
firewall-cmd --list-services &
N/A &
Lista los servicios permitidos. &
Linux
\\
\hline

firewall-cmd &
firewall-cmd --add-port &
--remove-port &
Permite abrir o cerrar puertos temporalmente o de manera permanente. &
Linux
\\
\hline

firewall-cmd &
firewall-cmd --reload &
N/A &
Recarga toda la configuración del firewall. &
Linux
\\
\hline

\end{longtable}

%%%%%%%%%%%%%%%%%%%%%%%%%%%%%%%%%%%%%%%%%%%%%%%%%%%%%%%%%%%%%%%%%%%%%%%%%%%%%%%

\subsubsection{Firewall de Windows}

Windows incorpora su propio firewall con una amplia interfaz de línea de
comandos disponible tanto desde CMD como desde PowerShell.

\begin{longtable}{|p{2.8cm}|p{3cm}|p{4cm}|p{5.6cm}|p{1.5cm}|}

\hline
Comando &
Sintaxis &
Argumentos útiles &
Particularidad &
OS
\\
\hline
\endhead

netsh &
netsh advfirewall show allprofiles &
N/A &
Consulta el estado del firewall para cada perfil de red. &
Windows
\\
\hline

netsh &
netsh advfirewall firewall show rule name=all &
N/A &
Lista todas las reglas configuradas. &
Windows
\\
\hline

netsh &
netsh advfirewall firewall add rule &
delete rule &
Permite crear o eliminar reglas desde CMD. &
Windows
\\
\hline

Get-NetFirewallRule &
Get-NetFirewallRule &
DisplayName &
Obtiene todas las reglas del firewall mediante PowerShell. &
Windows
\\
\hline

Get-NetFirewallPortFilter &
Get-NetFirewallPortFilter &
N/A &
Consulta los puertos asociados a una regla determinada. &
Windows
\\
\hline

Enable-NetFirewallRule &
Enable-NetFirewallRule &
Disable-NetFirewallRule &
Habilita o deshabilita reglas existentes. &
Windows
\\
\hline

Set-NetFirewallProfile &
Set-NetFirewallProfile &
Enabled &
Permite habilitar o deshabilitar perfiles completos del firewall. &
Windows
\\
\hline

\end{longtable}

%%%%%%%%%%%%%%%%%%%%%%%%%%%%%%%%%%%%%%%%%%%%%%%%%%%%%%%%%%%%%%%%%%%%%%%%%%%%%%%

\subsection{Paso 9. Captura y análisis de tráfico}

Cuando la configuración parece correcta y los servicios continúan sin responder
adecuadamente, resulta necesario inspeccionar directamente el tráfico de red.

Durante esta etapa se busca responder preguntas como:

\begin{itemize}
\item ¿Los paquetes realmente salen del equipo?
\item ¿Llegan respuestas?
\item ¿Existe retransmisión TCP?
\item ¿Hay pérdida de paquetes?
\item ¿El handshake TCP se completa?
\item ¿Existe tráfico DNS?
\item ¿La aplicación realmente está transmitiendo información?
\end{itemize}

%%%%%%%%%%%%%%%%%%%%%%%%%%%%%%%%%%%%%%%%%%%%%%%%%%%%%%%%%%%%%%%%%%%%%%%%%%%%%%%

\subsubsection{Captura básica de tráfico}

Las siguientes herramientas permiten capturar paquetes directamente desde una
interfaz de red.

\begin{longtable}{|p{2.8cm}|p{3cm}|p{4cm}|p{5.6cm}|p{1.5cm}|}

\hline
Comando &
Sintaxis &
Argumentos útiles &
Particularidad &
OS
\\
\hline
\endhead

tcpdump &
tcpdump -i eth0 &
-n, -vv, -c &
Herramienta estándar para capturar tráfico de red. Constituye una de las utilidades más importantes durante troubleshooting avanzado. &
Linux/macOS
\\
\hline

tcpdump &
tcpdump -D &
N/A &
Lista todas las interfaces disponibles para captura. &
Linux/macOS
\\
\hline

tcpdump &
tcpdump -i any &
N/A &
Captura tráfico proveniente de todas las interfaces disponibles. &
Linux
\\
\hline

tcpdump &
tcpdump -nn &
-v, -vv, -vvv &
Evita la resolución de nombres DNS y puertos, mostrando únicamente direcciones IP y números de puerto. Reduce el ruido durante la captura. &
Linux/macOS
\\
\hline

tcpdump &
tcpdump -c 100 &
-c &
Finaliza automáticamente la captura después de obtener el número indicado de paquetes. &
Linux/macOS
\\
\hline

tcpdump &
tcpdump -w captura.pcap &
-U &
Guarda la captura en formato PCAP para su posterior análisis. &
Linux/macOS
\\
\hline

tcpdump &
tcpdump -r captura.pcap &
-n &
Permite analizar una captura previamente almacenada. &
Linux/macOS
\\
\hline

\end{longtable}

%%%%5
%%%%%%%%%%%%%%%%%%%%%%%%%%%%%%%%%%%%%%%%%%%%%%%%%%%%%%%%%%%%%%%%%%%%%%%%%%%%%%%

\subsubsection{Filtros por host, puerto y protocolo}

Uno de los aspectos más importantes de tcpdump consiste en limitar la captura
únicamente al tráfico relevante.

\begin{longtable}{|p{2.8cm}|p{4.2cm}|p{3.5cm}|p{5cm}|p{1.5cm}|}

\hline
Comando &
Sintaxis &
Argumentos útiles &
Particularidad &
OS
\\
\hline
\endhead

tcpdump &
tcpdump host 192.168.1.15 &
src, dst &
Captura únicamente paquetes cuyo origen o destino corresponde al host indicado. &
Linux/macOS
\\
\hline

tcpdump &
tcpdump port 443 &
src, dst &
Captura tráfico asociado a un puerto específico. Muy utilizado para HTTPS, DNS, SSH, etc. &
Linux/macOS
\\
\hline

tcpdump &
tcpdump tcp &
udp, icmp &
Filtra únicamente un protocolo determinado. &
Linux/macOS
\\
\hline

tcpdump &
tcpdump src host 10.0.0.5 &
dst host &
Filtra únicamente paquetes provenientes de un origen específico. &
Linux/macOS
\\
\hline

tcpdump &
tcpdump dst port 443 &
src port &
Permite filtrar tráfico dirigido hacia un puerto específico. &
Linux/macOS
\\
\hline

tcpdump &
tcpdump 'tcp port 80 or tcp port 443' &
and, or, not &
Permite construir filtros booleanos complejos utilizando la sintaxis BPF. &
Linux/macOS
\\
\hline

tcpdump &
tcpdump net 192.168.10.0/24 &
mask &
Captura todo el tráfico correspondiente a una red determinada. &
Linux/macOS
\\
\hline

tcpdump &
tcpdump broadcast &
multicast &
Captura únicamente tráfico broadcast o multicast. &
Linux/macOS
\\
\hline

\end{longtable}


%%%%%%%%%%%5
%%%%%%%%%%%%%%%%%%%%%%%%%%%%%%%%%%%%%%%%%%%%%%%%%%%%%%%%%%%%%%%%%%%%%%%%%%%%%%%

\subsubsection{Capturas específicas para troubleshooting}

Las siguientes combinaciones constituyen algunas de las capturas más utilizadas
durante el diagnóstico de incidencias.

\begin{longtable}{|p{3cm}|p{4cm}|p{3cm}|p{5cm}|p{1.5cm}|}

\hline
Objetivo &
Comando &
Filtro &
Particularidad &
OS
\\
\hline
\endhead

DNS &
tcpdump port 53 &
udp &
Permite verificar consultas y respuestas DNS. &
Linux/macOS
\\
\hline

HTTPS &
tcpdump port 443 &
tcp &
Permite comprobar el establecimiento del handshake TCP previo a TLS. &
Linux/macOS
\\
\hline

HTTP &
tcpdump port 80 &
tcp &
Captura tráfico HTTP sin cifrar. &
Linux/macOS
\\
\hline

SSH &
tcpdump port 22 &
tcp &
Verifica conexiones SSH entrantes y salientes. &
Linux/macOS
\\
\hline

ICMP &
tcpdump icmp &
N/A &
Permite analizar solicitudes y respuestas Ping. &
Linux/macOS
\\
\hline

ARP &
tcpdump arp &
N/A &
Permite diagnosticar problemas de resolución MAC. &
Linux/macOS
\\
\hline

DHCP &
tcpdump port 67 or port 68 &
udp &
Captura solicitudes y respuestas DHCP. &
Linux/macOS
\\
\hline

NTP &
tcpdump port 123 &
udp &
Permite verificar sincronización horaria mediante NTP. &
Linux/macOS
\\
\hline

\end{longtable}
%%%%%%%%%%%%%%%%%%%%%%5

%%%%%%%%%%%%%%%%%%%%%%%%%%%%%%%%%%%%%%%%%%%%%%%%%%%%%%%%%%%%%%%%%%%%%%%%%%%%%%%

\subsubsection{Herramientas complementarias}

Existen otras herramientas orientadas al análisis de tráfico desde la línea de
comandos que resultan especialmente útiles durante tareas de troubleshooting.

\begin{longtable}{|p{2.8cm}|p{3cm}|p{4cm}|p{5.6cm}|p{1.5cm}|}

\hline
Comando &
Sintaxis &
Argumentos útiles &
Particularidad &
OS
\\
\hline
\endhead

tshark &
tshark -i eth0 &
-f, -Y, -w, -r &
Versión de línea de comandos de Wireshark. Permite capturar y analizar tráfico utilizando prácticamente las mismas capacidades del entorno gráfico. &
Linux/macOS/Windows
\\
\hline

tshark &
tshark -r captura.pcap &
-Y &
Analiza archivos PCAP utilizando filtros de visualización de Wireshark. &
Todos
\\
\hline

ngrep &
ngrep patrón &
port &
Permite buscar cadenas de texto dentro del tráfico de red. Muy útil para protocolos sin cifrar como HTTP, SMTP o FTP. &
Linux/macOS
\\
\hline

tcpflow &
tcpflow &
-c &
Reconstruye flujos TCP completos, facilitando el análisis de conversaciones entre cliente y servidor. &
Linux/macOS
\\
\hline

capinfos &
capinfos captura.pcap &
N/A &
Obtiene estadísticas generales de un archivo PCAP: duración, número de paquetes, tamaño y protocolos observados. &
Linux/macOS
\\
\hline

editcap &
editcap entrada.pcap salida.pcap &
-c, -A, -B &
Permite dividir, recortar o filtrar archivos PCAP antes de analizarlos. &
Linux/macOS
\\
\hline

mergecap &
mergecap &
-w &
Fusiona múltiples capturas PCAP en un único archivo. Muy útil cuando se captura tráfico desde varias interfaces o equipos. &
Linux/macOS
\\
\hline

\end{longtable}


%%%%%%%%%%%%%%%%%%%%%%%%%%%%%%%%%%%%%%%%%%%%%%%%%%%%%%%%%%%%%%%%%%%%%%%%%%%%%%%

\subsection{Paso 10. Revisión de logs y monitoreo del sistema}

Una vez verificada la conectividad y el estado de los servicios, el siguiente
paso consiste en revisar los registros generados por el sistema operativo y las
aplicaciones.

Durante esta etapa se busca responder preguntas como:

\begin{itemize}
\item ¿Qué error ocurrió?
\item ¿Cuándo ocurrió?
\item ¿Qué proceso lo generó?
\item ¿Existe algún patrón repetitivo?
\item ¿Qué archivo contiene la información relevante?
\item ¿El problema continúa ocurriendo en este momento?
\end{itemize}

%%%%%%%%%%%%%%%%%%%%%%%%%%%%%%%%%%%%%%%%%%%%%%%%%%%%%%%%%%%%%%%%%%%%%%%%%%%%%%%

\subsubsection{Visualización de logs}

Los siguientes comandos permiten inspeccionar archivos de registro.

\begin{longtable}{|p{2.8cm}|p{3.3cm}|p{4cm}|p{5.3cm}|p{1.5cm}|}

\hline
Comando &
Sintaxis &
Argumentos útiles &
Particularidad &
OS
\\
\hline
\endhead

cat &
cat archivo.log &
N/A &
Muestra completamente el contenido de un archivo de registro. Adecuado únicamente para archivos pequeños. &
Linux/macOS
\\
\hline

less &
less archivo.log &
-N, +G &
Permite navegar archivos grandes de forma interactiva sin cargarlos completamente en memoria. &
Linux/macOS
\\
\hline

more &
more archivo.log &
N/A &
Visualiza archivos página por página. &
Todos
\\
\hline

head &
head archivo.log &
-n &
Muestra las primeras líneas del archivo. Muy útil para revisar encabezados o información inicial. &
Linux/macOS
\\
\hline

tail &
tail archivo.log &
-n &
Muestra las últimas líneas del archivo. Normalmente es el primer comando utilizado durante troubleshooting. &
Linux/macOS
\\
\hline

tail &
tail -f archivo.log &
-f &
Permite seguir un log en tiempo real conforme se escriben nuevas entradas. &
Linux/macOS
\\
\hline

tail &
tail -F archivo.log &
-F &
Continúa siguiendo el archivo incluso si éste es rotado o recreado por logrotate. &
Linux/macOS
\\
\hline

type &
type archivo.log &
N/A &
Visualiza archivos desde CMD. &
Windows
\\
\hline

Get-Content &
Get-Content archivo.log &
-Tail, -Wait &
Equivalente a tail en PowerShell. El argumento -Wait permite seguir el archivo en tiempo real. &
Windows
\\
\hline

\end{longtable}

%%%%%%%%%%%%%%%%%%%%%%%%%%%%%%%%%%%%%%%%%%%%%%%%%%%%%%%%%%%%%%%%%%%%%%%%%%%%%%%

\subsubsection{Búsqueda dentro de logs}

Una vez localizado el archivo correcto, normalmente resulta necesario buscar
mensajes específicos.

\begin{longtable}{|p{2.8cm}|p{3.3cm}|p{4cm}|p{5.3cm}|p{1.5cm}|}

\hline
Comando &
Sintaxis &
Argumentos útiles &
Particularidad &
OS
\\
\hline
\endhead

grep &
grep ERROR archivo.log &
-i, -n, -r, -A, -B, -C, -v &
Busca cadenas o expresiones regulares dentro de archivos. Constituye una de las herramientas más utilizadas durante troubleshooting. &
Linux/macOS
\\
\hline

grep &
grep -i exception &
-E &
Permite realizar búsquedas ignorando mayúsculas/minúsculas o utilizando expresiones regulares extendidas. &
Linux/macOS
\\
\hline

findstr &
findstr ERROR archivo.log &
/I, /N, /S &
Equivalente básico de grep en Windows. &
Windows
\\
\hline

awk &
awk &
-F &
Permite analizar columnas y generar reportes a partir de archivos de log. &
Linux/macOS
\\
\hline

sed &
sed &
-n &
Filtra o transforma líneas específicas del archivo. &
Linux/macOS
\\
\hline

sort &
sort archivo.log &
-r, -u &
Ordena registros antes de analizarlos. &
Linux/macOS
\\
\hline

uniq &
uniq -c &
-d &
Cuenta ocurrencias repetidas de un mismo mensaje. Muy útil para detectar errores recurrentes. &
Linux/macOS
\\
\hline

wc &
wc -l archivo.log &
-l &
Cuenta rápidamente el número de líneas contenidas en un log. &
Linux/macOS
\\
\hline


\end{longtable}

%%%%%%%%%%%%%%%%%%%%%%%%%%%%%%%%%%%%%%%%%%%%%%%%%%%%%%%%%%%%%%%%%%%%%%%%%%%%%%%

\subsubsection{Logs del sistema}

Muchos errores no aparecen en los registros de la aplicación sino en los logs
del propio sistema operativo.

\begin{longtable}{|p{2.8cm}|p{3.3cm}|p{4cm}|p{5.3cm}|p{1.5cm}|}

\hline
Comando &
Sintaxis &
Argumentos útiles &
Particularidad &
OS
\\
\hline
\endhead

journalctl &
journalctl &
-x, -b, -f, -u &
Permite consultar todos los registros administrados por systemd. Constituye la principal herramienta de diagnóstico en distribuciones Linux modernas. &
Linux
\\
\hline

journalctl &
journalctl -u nginx &
-f &
Consulta únicamente los registros correspondientes a un servicio específico. &
Linux
\\
\hline

journalctl &
journalctl --since "10 min ago" &
--until &
Permite consultar únicamente eventos ocurridos dentro de un intervalo de tiempo. &
Linux
\\
\hline

journalctl &
journalctl -p err &
warning, info &
Filtra registros según su nivel de severidad. &
Linux
\\
\hline

journalctl &
journalctl -b &
-1 &
Consulta únicamente los registros del arranque actual o de arranques anteriores. &
Linux
\\
\hline

dmesg &
dmesg &
-T, -w &
Consulta mensajes generados por el kernel. Muy útil para diagnosticar problemas de hardware, controladores o dispositivos. &
Linux
\\
\hline

log show &
log show &
--last, --predicate &
Consulta el sistema unificado de logs de macOS. &
macOS
\\
\hline

wevtutil &
wevtutil qe System &
Application &
Consulta registros del Event Viewer desde la línea de comandos. &
Windows
\\
\hline

Get-WinEvent &
Get-WinEvent &
-LogName, -MaxEvents &
Permite consultar registros del Event Viewer utilizando PowerShell. &
Windows
\\
\hline

\end{longtable}

%%%%%%%%%%%%%%%%%%%%%%%%%%%%%%%%%%%%%%%%%%%%%%%%%%%%%%%%%%%%%%%%%%%%%%%%%%%%%%%

\subsubsection{Monitoreo en tiempo real}

Cuando un problema ocurre de manera intermitente, resulta útil observar el
sistema conforme los eventos suceden.

\begin{longtable}{|p{2.8cm}|p{3.3cm}|p{4cm}|p{5.3cm}|p{1.5cm}|}

\hline
Comando &
Sintaxis &
Argumentos útiles &
Particularidad &
OS
\\
\hline
\endhead

tail &
tail -f archivo.log &
-F &
Monitorea continuamente un archivo de registro. &
Linux/macOS
\\
\hline

journalctl &
journalctl -f &
-u &
Permite seguir en tiempo real los registros del sistema o de un servicio específico. &
Linux
\\
\hline

watch &
watch comando &
-n &
Ejecuta periódicamente un comando mostrando los cambios ocurridos entre ejecuciones. &
Linux
\\
\hline

Get-Content &
Get-Content archivo.log -Wait &
-Tail &
Permite seguir un archivo de log desde PowerShell. &
Windows
\\
\hline

\end{longtable}

%%%%%%%%%%%%%%%%%%%%%%%%%%%%%%%%%%%%%%%%%%%%%%%%%%%%%%%%%%%%%%%%%%%%%%%%%%%%%%%

\subsubsection{Comandos compuestos de troubleshooting}

Durante una incidencia es común combinar varias herramientas para localizar
rápidamente la causa del problema. Los siguientes comandos representan algunos
de los escenarios más frecuentes encontrados durante tareas de soporte.

\begin{longtable}{|p{4.8cm}|p{7.5cm}|p{4cm}|}

\hline
Objetivo &
Comando &
Particularidad
\\
\hline
\endhead

Buscar todos los archivos modificados durante la última hora &
\texttt{find / -type f -mmin -60 2>/dev/null} &
Muy útil para localizar archivos generados o modificados recientemente.
\\
\hline

Buscar archivos modificados hoy &
\texttt{find /var/log -type f -daystart -mtime 0} &
Permite reducir rápidamente la búsqueda de logs recientes.
\\
\hline

Buscar archivos mayores a 1 GB &
\texttt{find / -type f -size +1G 2>/dev/null} &
Ideal para investigar problemas de almacenamiento.
\\
\hline

Buscar archivos cuyo nombre contenga "log" &
\texttt{find / -iname "*log*" 2>/dev/null} &
Permite localizar rápidamente archivos de registro cuyo nombre se desconoce.
\\
\hline

Encontrar archivos que están creciendo constantemente &
\texttt{watch 'ls -lh /var/log'} &
Permite observar cambios de tamaño en tiempo real.
\\
\hline

Encontrar los archivos modificados más recientemente &
\texttt{find /var/log -type f -printf "\%T@ \%p$ \ $n" | sort -nr} &
Ordena los archivos según su última fecha de modificación.
\\
\hline

Visualizar simultáneamente varios logs &
\texttt{tail -f log1.log log2.log log3.log} &
Permite seguir múltiples registros desde una única terminal.
\\
\hline

Buscar únicamente errores &
\texttt{grep -i "error\textbackslash|exception\textbackslash|fatal" archivo.log} &
Filtra rápidamente los mensajes más relevantes.
\\
\hline

Contar cuántas veces aparece un error &
\texttt{grep ERROR archivo.log | wc -l} &
Permite medir la frecuencia de un problema.
\\
\hline

Encontrar los errores más repetidos &
\texttt{grep ERROR archivo.log | sort | uniq -c | sort -nr} &
Muy útil para detectar patrones repetitivos.
\\
\hline

Buscar registros generados durante los últimos diez minutos &
\texttt{journalctl --since "10 minutes ago"} &
Evita revisar grandes cantidades de información innecesaria.
\\
\hline

Seguir únicamente los logs de un servicio &
\texttt{journalctl -fu nginx} &
Monitorea un servicio mientras se reproduce el problema.
\\
\hline

Buscar un PID y verificar si continúa ejecutándose &
\texttt{ps -ef | grep PID} &
Permite comprobar rápidamente si un proceso terminó inesperadamente.
\\
\hline

Buscar procesos por nombre &
\texttt{ps -ef | grep nombreProceso} &
Método clásico para localizar aplicaciones en ejecución.
\\
\hline

Encontrar qué proceso utiliza un puerto &
\texttt{lsof -i :443} &
Muy útil cuando un servicio no puede iniciar por conflicto de puertos.
\\
\hline

Encontrar quién mantiene abierto un archivo &
\texttt{lsof archivo.log} &
Permite identificar procesos que bloquean archivos.
\\
\hline

Buscar procesos que escriben dentro de un directorio &
\texttt{lsof +D /var/log} &
Muestra todos los procesos utilizando archivos dentro del directorio.
\\
\hline

Liberar un archivo ocupado &
\texttt{fuser -k archivo.log} &
Finaliza el proceso que mantiene abierto el archivo.
\\
\hline

Encontrar archivos eliminados que continúan abiertos &
\texttt{lsof | grep deleted} &
Caso muy frecuente cuando el disco permanece lleno aun después de borrar archivos.
\\
\hline

Buscar todos los archivos abiertos por un proceso &
\texttt{lsof -p PID} &
Permite inspeccionar completamente la actividad de un proceso.
\\
\hline

Comprobar qué proceso consume más CPU &
\texttt{ps -eo pid,comm,\%cpu --sort=-\%cpu | head} &
Obtiene rápidamente los procesos con mayor utilización de CPU.
\\
\hline

Comprobar qué proceso consume más memoria &
\texttt{ps -eo pid,comm,\%mem --sort=-\%mem | head} &
Muy útil para detectar fugas de memoria.
\\
\hline

Buscar archivos cuyo tamaño cambia constantemente &
\texttt{watch 'du -sh /var/log/*'} &
Permite identificar qué archivo está creciendo.
\\
\hline

Buscar qué proceso escribe continuamente en disco &
\texttt{iotop} &
Muestra el uso de E/S por proceso en tiempo real.
\\
\hline

Encontrar archivos bloqueados &
\texttt{lsof | grep "(deleted)"} &
Frecuente cuando una aplicación continúa escribiendo en un archivo eliminado.
\\
\hline

Comprobar espacio disponible y uso por directorio &
\texttt{df -h ; du -sh /*} &
Permite determinar rápidamente dónde se encuentra el consumo de almacenamiento.
\\
\hline

Buscar conexiones hacia un puerto específico &
\texttt{ss -tan | grep :443} &
Permite verificar si realmente llegan conexiones al servidor.
\\
\hline

Verificar conexiones en estado LISTEN &
\texttt{ss -ltnp} &
Permite confirmar que el servicio realmente abrió el puerto esperado.
\\
\hline

Buscar retransmisiones TCP &
\texttt{netstat -s | grep retrans} &
Puede indicar pérdida de paquetes o problemas de red.
\\
\hline

Verificar rápidamente resolución DNS y conexión HTTPS &
\texttt{curl -Iv https://servidor.com} &
Permite comprobar DNS, TCP, TLS y HTTP mediante un único comando.
\\
\hline

\end{longtable}

%%%%%%%%%%%%%%%%%%%%%%%%%%%%%%%%%%%%%%%%%%%%%%%%%%%%%%%%%%%%%%%%%%%%%%%%%%%%%%%

\subsection{Paso 11. Administración de procesos y servicios}

Durante el proceso de troubleshooting resulta indispensable verificar que los
servicios se encuentren en ejecución, que no hayan terminado inesperadamente y
que consuman los recursos esperados.

Las siguientes herramientas permiten consultar, administrar y monitorear
procesos y servicios del sistema.

%%%%%%%%%%%%%%%%%%%%%%%%%%%%%%%%%%%%%%%%%%%%%%%%%%%%%%%%%%%%%%%%%%%%%%%%%%%%%%%

\subsubsection{Procesos}

\begin{longtable}{|p{2.8cm}|p{3.2cm}|p{4cm}|p{5.4cm}|p{1.5cm}|}

\hline
Comando &
Sintaxis &
Argumentos útiles &
Particularidad &
OS
\\
\hline
\endhead

ps &
ps -ef &
aux, -eo &
Lista todos los procesos en ejecución. Constituye el punto de partida para localizar procesos específicos. &
Linux/macOS
\\
\hline

ps &
ps -eo pid,user,comm,\%cpu,\%mem &
--sort &
Permite personalizar las columnas mostradas y ordenar los resultados según distintos criterios. &
Linux/macOS
\\
\hline

top &
top &
-p, -H &
Monitorea procesos en tiempo real mostrando utilización de CPU, memoria y carga del sistema. &
Linux/macOS
\\
\hline

htop &
htop &
-d &
Versión interactiva de top. Facilita ordenar procesos, buscar aplicaciones y finalizar procesos directamente desde la interfaz. &
Linux
\\
\hline

pidof &
pidof proceso &
N/A &
Obtiene rápidamente el PID asociado a un proceso. &
Linux
\\
\hline

pgrep &
pgrep nombre &
-a, -f &
Busca procesos utilizando expresiones regulares o parte del nombre del comando. &
Linux/macOS
\\
\hline

pkill &
pkill nombre &
-f &
Finaliza procesos utilizando su nombre en lugar del PID. &
Linux/macOS
\\
\hline

kill &
kill PID &
-9, -15 &
Envía señales a un proceso. SIGTERM permite finalizar ordenadamente; SIGKILL fuerza la terminación inmediata. &
Linux/macOS
\\
\hline

killall &
killall proceso &
-s &
Finaliza todos los procesos cuyo nombre coincida con el indicado. &
Linux/macOS
\\
\hline

nice &
nice -n 10 comando &
N/A &
Ejecuta un proceso con una prioridad determinada. &
Linux/macOS
\\
\hline

renice &
renice -10 PID &
N/A &
Modifica la prioridad de un proceso ya existente. &
Linux/macOS
\\
\hline

tasklist &
tasklist &
/FI, /SVC &
Lista todos los procesos en Windows junto con información adicional. &
Windows
\\
\hline

Get-Process &
Get-Process &
-Id, -Name &
Obtiene información detallada de procesos mediante PowerShell. &
Windows
\\
\hline

Stop-Process &
Stop-Process &
-Id, -Name, -Force &
Finaliza procesos desde PowerShell. &
Windows
\\
\hline

taskkill &
taskkill /PID PID &
/F, /IM &
Finaliza procesos utilizando el PID o el nombre del ejecutable. &
Windows
\\
\hline

\end{longtable}

%%%%%%%%%%%%%%%%%%%%%%%%%%%%%%%%%%%%%%%%%%%%%%%%%%%%%%%%%%%%%%%%%%%%%%%%%%%%%%%

\subsubsection{Servicios}

La mayoría de las aplicaciones empresariales se ejecutan como servicios del
sistema operativo. Las siguientes herramientas permiten consultar su estado y
administrarlos.

\begin{longtable}{|p{2.8cm}|p{3.2cm}|p{4cm}|p{5.4cm}|p{1.5cm}|}

\hline
Comando &
Sintaxis &
Argumentos útiles &
Particularidad &
OS
\\
\hline
\endhead

systemctl &
systemctl status servicio &
--no-pager &
Consulta el estado actual de un servicio administrado por systemd. Constituye normalmente el primer comando utilizado cuando una aplicación deja de funcionar. &
Linux
\\
\hline

systemctl &
systemctl start servicio &
stop, restart, reload &
Permite iniciar, detener o reiniciar servicios. &
Linux
\\
\hline

systemctl &
systemctl enable servicio &
disable &
Configura si el servicio debe iniciarse automáticamente durante el arranque. &
Linux
\\
\hline

systemctl &
systemctl list-units &
--type=service &
Lista todos los servicios administrados por systemd. &
Linux
\\
\hline

systemctl &
systemctl is-active servicio &
is-enabled &
Comprueba rápidamente si un servicio está activo o habilitado. &
Linux
\\
\hline

service &
service servicio status &
start, stop, restart &
Interfaz tradicional utilizada en distribuciones sin systemd o por compatibilidad. &
Linux
\\
\hline

sc &
sc query &
queryex &
Consulta servicios desde la línea de comandos de Windows. &
Windows
\\
\hline

sc &
sc start servicio &
stop &
Permite iniciar o detener servicios de Windows. &
Windows
\\
\hline

Get-Service &
Get-Service &
-Name &
Obtiene información sobre servicios mediante PowerShell. &
Windows
\\
\hline

Restart-Service &
Restart-Service &
-Name &
Reinicia servicios desde PowerShell. &
Windows
\\
\hline

net &
net start &
stop &
Herramienta clásica para iniciar y detener servicios en Windows. &
Windows
\\
\hline

launchctl &
launchctl list &
print &
Permite consultar los servicios y demonios administrados por launchd en macOS. &
macOS
\\
\hline

launchctl &
launchctl load &
unload &
Carga o descarga servicios en macOS. &
macOS
\\
\hline

\end{longtable}

%%%%%%%%%%%%%%%%%%%%%%%%%%%%%%%%%%%%%%%%%%%%%%%%%%%%%%%%%%%%%%%%%%%%%%%%%%%%%%%

\subsection{Paso 12. Verificación de recursos del sistema}

Una vez comprobado que los servicios se encuentran en ejecución, resulta
conveniente verificar el estado general del sistema operativo y la
disponibilidad de sus recursos.

Durante esta etapa se busca responder preguntas como:

\begin{itemize}
\item ¿Existe saturación de CPU?
\item ¿Hay memoria RAM disponible?
\item ¿Se está utilizando memoria swap?
\item ¿Existe espera por operaciones de disco?
\item ¿La carga del sistema es anormal?
\item ¿Hay procesos consumiendo recursos excesivos?
\end{itemize}

%%%%%%%%%%%%%%%%%%%%%%%%%%%%%%%%%%%%%%%%%%%%%%%%%%%%%%%%%%%%%%%%%%%%%%%%%%%%%%%

\subsubsection{CPU y carga del sistema}

\begin{longtable}{|p{2.8cm}|p{3.2cm}|p{4cm}|p{5.4cm}|p{1.5cm}|}

\hline
Comando &
Sintaxis &
Argumentos útiles &
Particularidad &
OS
\\
\hline
\endhead

uptime &
uptime &
N/A &
Muestra el tiempo desde el último arranque y el Load Average del sistema. Es uno de los primeros comandos utilizados para evaluar la carga general. &
Linux/macOS
\\
\hline

top &
top &
-d, -H &
Permite observar en tiempo real el consumo de CPU por proceso y la carga global del sistema. &
Linux/macOS
\\
\hline

htop &
htop &
-d &
Versión interactiva de top con ordenamiento, búsqueda y filtrado de procesos. &
Linux
\\
\hline

mpstat &
mpstat &
-P ALL &
Muestra estadísticas de utilización para cada núcleo del procesador. Muy útil para detectar cargas desbalanceadas. &
Linux
\\
\hline

sar &
sar -u &
-u, -P, -q &
Permite consultar estadísticas históricas de utilización de CPU cuando el servicio sysstat se encuentra habilitado. &
Linux
\\
\hline

vmstat &
vmstat 1 &
-s &
Resume el estado general del sistema incluyendo CPU, memoria y procesos. &
Linux/macOS
\\
\hline

wmic &
wmic cpu get loadpercentage &
N/A &
Obtiene el porcentaje de utilización del procesador desde CMD. &
Windows
\\
\hline

Get-Counter &
Get-Counter "\textbackslash Processor(\_Total)\textbackslash \% Processor Time" &
N/A &
Consulta contadores de rendimiento de CPU mediante PowerShell. &
Windows
\\
\hline

\end{longtable}

%%%%%%%%%%%%%%%%%%%%%%%%%%%%%%%%%%%%%%%%%%%%%%%%%%%%%%%%%%%%%%%%%%%%%%%%%%%%%%%

\subsubsection{Memoria}

\begin{longtable}{|p{2.8cm}|p{3.2cm}|p{4cm}|p{5.4cm}|p{1.5cm}|}

\hline
Comando &
Sintaxis &
Argumentos útiles &
Particularidad &
OS
\\
\hline
\endhead

free &
free -h &
-m, -g &
Muestra el uso de memoria RAM y swap en formatos legibles. Constituye el comando más utilizado para comprobar disponibilidad de memoria. &
Linux
\\
\hline

vmstat &
vmstat &
1 &
Permite observar el uso de memoria, swap y actividad del sistema en tiempo real. &
Linux/macOS
\\
\hline

cat &
cat /proc/meminfo &
N/A &
Presenta información detallada sobre la memoria del sistema. &
Linux
\\
\hline

swapon &
swapon --show &
-s &
Lista las áreas de swap configuradas y su utilización. &
Linux
\\
\hline

vm\_stat &
vm\_stat &
N/A &
Consulta estadísticas de memoria virtual en macOS. &
macOS
\\
\hline

systeminfo &
systeminfo &
N/A &
Incluye información general de memoria física instalada y disponible. &
Windows
\\
\hline

Get-ComputerInfo &
Get-ComputerInfo &
N/A &
Obtiene información completa del sistema mediante PowerShell, incluyendo memoria instalada. &
Windows
\\
\hline

Get-Counter &
Get-Counter "\textbackslash Memory\textbackslash Available MBytes" &
N/A &
Consulta la memoria disponible utilizando contadores de rendimiento. &
Windows
\\
\hline

\end{longtable}

%%%%%%%%%%%%%%%%%%%%%%%%%%%%%%%%%%%%%%%%%%%%%%%%%%%%%%%%%%%%%%%%%%%%%%%%%%%%%%%

\subsubsection{Disco e I/O}

\begin{longtable}{|p{2.8cm}|p{3.2cm}|p{4cm}|p{5.4cm}|p{1.5cm}|}

\hline
Comando &
Sintaxis &
Argumentos útiles &
Particularidad &
OS
\\
\hline
\endhead

iostat &
iostat &
-x, -d &
Presenta estadísticas de entrada y salida por dispositivo de almacenamiento. Muy útil para detectar cuellos de botella de disco. &
Linux
\\
\hline

iotop &
iotop &
-o &
Muestra qué procesos están realizando operaciones intensivas de lectura o escritura en disco. &
Linux
\\
\hline

sar &
sar -d &
-d &
Obtiene estadísticas históricas de utilización de dispositivos de almacenamiento. &
Linux
\\
\hline

dstat &
dstat &
-c, -d, -m, -n &
Combina estadísticas de CPU, disco, memoria y red en una única vista. &
Linux
\\
\hline

vmstat &
vmstat 1 &
N/A &
Permite observar bloqueos relacionados con operaciones de entrada y salida. &
Linux/macOS
\\
\hline

Get-PhysicalDisk &
Get-PhysicalDisk &
N/A &
Consulta información sobre discos físicos desde PowerShell. &
Windows
\\
\hline

Get-Counter &
Get-Counter "\textbackslash PhysicalDisk(*)\textbackslash Disk Transfers/sec" &
N/A &
Consulta estadísticas de actividad de disco. &
Windows
\\
\hline

\end{longtable}

%%%%%%%%%%%%%%%%%%%%%%%%%%%%%%%%%%%%%%%%%%%%%%%%%%%%%%%%%%%%%%%%%%%%%%%%%%%%%%%

\subsubsection{Información general del sistema}

Además de los recursos, suele ser útil obtener información general del sistema
operativo durante una incidencia.

\begin{longtable}{|p{2.8cm}|p{3.2cm}|p{4cm}|p{5.4cm}|p{1.5cm}|}

\hline
Comando &
Sintaxis &
Argumentos útiles &
Particularidad &
OS
\\
\hline
\endhead

uname &
uname -a &
-r, -m &
Muestra información del kernel, arquitectura y sistema operativo. &
Linux/macOS
\\
\hline

hostnamectl &
hostnamectl &
status &
Presenta información del host, sistema operativo y kernel en sistemas con systemd. &
Linux
\\
\hline

hostname &
hostname &
-f, -I &
Consulta el nombre del equipo y, según el sistema, direcciones IP asociadas. &
Linux/macOS
\\
\hline

lscpu &
lscpu &
N/A &
Obtiene información detallada del procesador. &
Linux
\\
\hline

lsmem &
lsmem &
N/A &
Presenta la distribución de la memoria física. &
Linux
\\
\hline

lsblk &
lsblk &
-f &
Lista dispositivos de almacenamiento, particiones y sistemas de archivos. &
Linux
\\
\hline

systeminfo &
systeminfo &
N/A &
Obtiene un resumen completo del sistema operativo Windows. &
Windows
\\
\hline

Get-ComputerInfo &
Get-ComputerInfo &
N/A &
Equivalente en PowerShell para consultar información del sistema. &
Windows
\\
\hline

sw\_vers &
sw\_vers &
N/A &
Consulta la versión del sistema operativo macOS. &
macOS
\\
\hline

\end{longtable}

%%%%%%%%%%%%%%%%%%%%%%%%%%%%%%%%%%%%%%%%%%%%%%%%%%%%%%%%%%%%%%%%%%%%%%%%%%%%%%%

\subsection{Paso 13. Verificación del sistema de archivos y almacenamiento}

Una parte importante del proceso de troubleshooting consiste en verificar el
estado del almacenamiento, los sistemas de archivos y los puntos de montaje.

Durante esta etapa se busca responder preguntas como:

\begin{itemize}
\item ¿Existe espacio libre en disco?
\item ¿El sistema de archivos está montado?
\item ¿Se encuentra en modo de solo lectura?
\item ¿Se agotaron los inodos?
\item ¿El dispositivo fue detectado correctamente?
\item ¿Los permisos del archivo son correctos?
\end{itemize}

%%%%%%%%%%%%%%%%%%%%%%%%%%%%%%%%%%%%%%%%%%%%%%%%%%%%%%%%%%%%%%%%%%%%%%%%%%%%%%%

\subsubsection{Espacio en disco}

\begin{longtable}{|p{2.8cm}|p{3.2cm}|p{4cm}|p{5.4cm}|p{1.5cm}|}

\hline
Comando &
Sintaxis &
Argumentos útiles &
Particularidad &
OS
\\
\hline
\endhead

df &
df -h &
-T, -i &
Muestra el espacio utilizado y disponible por sistema de archivos. Constituye el primer comando utilizado para diagnosticar problemas de almacenamiento. &
Linux/macOS
\\
\hline

df &
df -i &
-h &
Permite comprobar el consumo de inodos. Un sistema puede quedarse sin inodos aun teniendo espacio libre. &
Linux
\\
\hline

du &
du -sh directorio &
-h, -d &
Calcula el espacio ocupado por archivos o directorios específicos. Muy útil para localizar consumo excesivo de almacenamiento. &
Linux/macOS
\\
\hline

du &
du -sh * &
--max-depth &
Resume el tamaño ocupado por cada directorio del nivel actual. &
Linux/macOS
\\
\hline

ncdu &
ncdu &
N/A &
Herramienta interactiva para explorar el uso del almacenamiento por directorios. Muy útil durante troubleshooting. &
Linux
\\
\hline

Get-PSDrive &
Get-PSDrive &
N/A &
Consulta la utilización de las unidades de almacenamiento desde PowerShell. &
Windows
\\
\hline

wmic &
wmic logicaldisk get size,freespace,caption &
N/A &
Obtiene el espacio libre de las unidades en Windows. &
Windows
\\
\hline

\end{longtable}

%%%%%%%%%%%%%%%%%%%%%%%%%%%%%%%%%%%%%%%%%%%%%%%%%%%%%%%%%%%%%%%%%%%%%%%%%%%%%%%

\subsubsection{Dispositivos y particiones}

\begin{longtable}{|p{2.8cm}|p{3.2cm}|p{4cm}|p{5.4cm}|p{1.5cm}|}

\hline
Comando &
Sintaxis &
Argumentos útiles &
Particularidad &
OS
\\
\hline
\endhead

lsblk &
lsblk &
-f, -o &
Lista discos, particiones, UUID y puntos de montaje. Constituye normalmente el primer comando para verificar almacenamiento. &
Linux
\\
\hline

blkid &
blkid &
N/A &
Obtiene UUID, etiquetas y tipos de sistema de archivos. &
Linux
\\
\hline

fdisk &
fdisk -l &
N/A &
Lista las particiones detectadas por el sistema. &
Linux
\\
\hline

parted &
parted -l &
N/A &
Consulta la tabla de particiones de todos los discos disponibles. &
Linux
\\
\hline

mount &
mount &
-t &
Lista todos los sistemas de archivos actualmente montados. &
Linux/macOS
\\
\hline

findmnt &
findmnt &
-t &
Presenta una vista jerárquica de los sistemas de archivos montados. &
Linux
\\
\hline

diskutil &
diskutil list &
info &
Herramienta principal para administrar discos en macOS. &
macOS
\\
\hline

Get-Disk &
Get-Disk &
N/A &
Lista discos físicos disponibles mediante PowerShell. &
Windows
\\
\hline

Get-Volume &
Get-Volume &
N/A &
Consulta volúmenes y sistemas de archivos montados en Windows. &
Windows
\\
\hline

\end{longtable}

%%%%%%%%%%%%%%%%%%%%%%%%%%%%%%%%%%%%%%%%%%%%%%%%%%%%%%%%%%%%%%%%%%%%%%%%%%%%%%%

\subsubsection{Montaje de sistemas de archivos}

\begin{longtable}{|p{2.8cm}|p{3.2cm}|p{4cm}|p{5.4cm}|p{1.5cm}|}

\hline
Comando &
Sintaxis &
Argumentos útiles &
Particularidad &
OS
\\
\hline
\endhead

mount &
mount dispositivo directorio &
-o, -t &
Permite montar manualmente un sistema de archivos. &
Linux/macOS
\\
\hline

umount &
umount directorio &
-l, -f &
Desmonta un sistema de archivos. Muy útil antes de realizar mantenimiento o retirar dispositivos. &
Linux/macOS
\\
\hline

mountpoint &
mountpoint directorio &
-q &
Permite comprobar si un directorio constituye un punto de montaje. &
Linux
\\
\hline

findmnt &
findmnt directorio &
N/A &
Obtiene información específica sobre un punto de montaje. &
Linux
\\
\hline

fsck &
fsck dispositivo &
-y, -f &
Verifica y repara inconsistencias en sistemas de archivos. Debe ejecutarse preferentemente sobre sistemas desmontados. &
Linux
\\
\hline

e2fsck &
e2fsck dispositivo &
-f &
Versión especializada para sistemas de archivos ext2/ext3/ext4. &
Linux
\\
\hline

xfs\_repair &
xfs\_repair dispositivo &
-L &
Repara sistemas de archivos XFS. &
Linux
\\
\hline

diskutil &
diskutil verifyVolume &
repairVolume &
Verifica y repara volúmenes en macOS. &
macOS
\\
\hline

chkdsk &
chkdsk unidad: &
/f, /r &
Herramienta clásica para comprobar y reparar sistemas de archivos NTFS/FAT. &
Windows
\\
\hline

\end{longtable}

%%%%%%%%%%%%%%%%%%%%%%%%%%%%%%%%%%%%%%%%%%%%%%%%%%%%%%%%%%%%%%%%%%%%%%%%%%%%%%%

\subsubsection{Información de archivos}

\begin{longtable}{|p{2.8cm}|p{3.2cm}|p{4cm}|p{5.4cm}|p{1.5cm}|}

\hline
Comando &
Sintaxis &
Argumentos útiles &
Particularidad &
OS
\\
\hline
\endhead

stat &
stat archivo &
-c &
Presenta información detallada sobre un archivo: permisos, propietario, tamaño, fechas e inodo. Es una de las herramientas más útiles para troubleshooting relacionado con archivos. &
Linux/macOS
\\
\hline

file &
file archivo &
-i &
Determina el tipo real de un archivo analizando su contenido y no únicamente su extensión. &
Linux/macOS
\\
\hline

readlink &
readlink enlace &
-f &
Resuelve enlaces simbólicos mostrando la ruta real del archivo. &
Linux/macOS
\\
\hline

realpath &
realpath archivo &
N/A &
Obtiene la ruta absoluta y resuelve enlaces simbólicos. &
Linux/macOS
\\
\hline

ls &
ls -li &
-l, -a, -h &
Permite consultar permisos, propietarios e inodos de archivos y directorios. &
Linux/macOS
\\
\hline

attrib &
attrib archivo &
+s, -s &
Consulta y modifica atributos de archivos en Windows. &
Windows
\\
\hline

Get-Item &
Get-Item archivo &
Force &
Obtiene propiedades detalladas de un archivo mediante PowerShell. &
Windows
\\
\hline

\end{longtable}


%%%%%%%%%%%%%%%%%%%%%%%%%%%%%%%%%%%%%%%%%%%%%%%%%%%%%%%%%%%%%%%%%%%%%%%%%%%%%%%

\subsection{Paso 14. Verificación de usuarios, grupos y permisos}

Una causa frecuente de incidencias corresponde a permisos insuficientes sobre
archivos, directorios o recursos del sistema.

Durante esta etapa se busca responder preguntas como:

\begin{itemize}
\item ¿Con qué usuario se está ejecutando la aplicación?
\item ¿Quién es el propietario del archivo?
\item ¿Qué permisos posee?
\item ¿El usuario pertenece al grupo correcto?
\item ¿Existen ACLs adicionales?
\item ¿SELinux o AppArmor están restringiendo el acceso?
\end{itemize}

%%%%%%%%%%%%%%%%%%%%%%%%%%%%%%%%%%%%%%%%%%%%%%%%%%%%%%%%%%%%%%%%%%%%%%%%%%%%%%%

\subsubsection{Usuarios y grupos}

\begin{longtable}{|p{2.8cm}|p{3.2cm}|p{4cm}|p{5.4cm}|p{1.5cm}|}

\hline
Comando &
Sintaxis &
Argumentos útiles &
Particularidad &
OS
\\
\hline
\endhead

whoami &
whoami &
N/A &
Muestra el usuario con el que se está ejecutando la sesión actual. Constituye normalmente la primera verificación relacionada con permisos. &
Linux/macOS
\\
\hline

id &
id &
-u, -g &
Obtiene el UID, GID y grupos secundarios del usuario actual o de un usuario específico. &
Linux/macOS
\\
\hline

groups &
groups usuario &
N/A &
Lista los grupos a los que pertenece un usuario. &
Linux/macOS
\\
\hline

getent &
getent passwd usuario &
group &
Consulta usuarios y grupos utilizando las bases de datos configuradas por el sistema (locales, LDAP, etc.). &
Linux
\\
\hline

who &
who &
-a &
Lista los usuarios conectados al sistema. &
Linux/macOS
\\
\hline

w &
w &
-h &
Muestra usuarios conectados junto con la actividad que realizan. &
Linux
\\
\hline

users &
users &
N/A &
Presenta únicamente los nombres de los usuarios conectados. &
Linux/macOS
\\
\hline

logname &
logname &
N/A &
Obtiene el usuario que inició la sesión actual. &
Linux/macOS
\\
\hline

query user &
query user &
N/A &
Lista las sesiones activas en Windows. &
Windows
\\
\hline

whoami &
whoami /all &
/groups, /priv &
Muestra información completa del usuario, SID, privilegios y grupos de seguridad. Muy útil durante troubleshooting de permisos en Windows. &
Windows
\\
\hline

Get-LocalUser &
Get-LocalUser &
Name &
Obtiene usuarios locales mediante PowerShell. &
Windows
\\
\hline

Get-LocalGroup &
Get-LocalGroup &
N/A &
Lista los grupos locales existentes. &
Windows
\\
\hline

\end{longtable}

%%%%%%%%%%%%%%%%%%%%%%%%%%%%%%%%%%%%%%%%%%%%%%%%%%%%%%%%%%%%%%%%%%%%%%%%%%%%%%%

\subsubsection{Permisos y propietarios}

\begin{longtable}{|p{2.8cm}|p{3.2cm}|p{4cm}|p{5.4cm}|p{1.5cm}|}

\hline
Comando &
Sintaxis &
Argumentos útiles &
Particularidad &
OS
\\
\hline
\endhead

ls &
ls -l &
-a, -h, -i &
Permite visualizar propietario, grupo y permisos POSIX de archivos y directorios. &
Linux/macOS
\\
\hline

stat &
stat archivo &
-c &
Presenta información detallada sobre permisos, propietario, fechas e inodo. &
Linux/macOS
\\
\hline

chmod &
chmod 755 archivo &
-R &
Modifica permisos POSIX utilizando notación octal o simbólica. &
Linux/macOS
\\
\hline

chown &
chown usuario archivo &
-R &
Cambia el propietario de archivos y directorios. &
Linux/macOS
\\
\hline

chgrp &
chgrp grupo archivo &
-R &
Modifica el grupo propietario de archivos y directorios. &
Linux/macOS
\\
\hline

umask &
umask &
-S &
Consulta o modifica la máscara de permisos por defecto utilizada al crear archivos nuevos. &
Linux/macOS
\\
\hline

icacls &
icacls archivo &
/grant, /deny, /inheritance &
Consulta y modifica listas de control de acceso (ACL) en Windows. &
Windows
\\
\hline

takeown &
takeown /F archivo &
/R &
Permite tomar posesión de archivos o directorios. &
Windows
\\
\hline

Get-Acl &
Get-Acl archivo &
N/A &
Obtiene las ACL de un archivo mediante PowerShell. &
Windows
\\
\hline

Set-Acl &
Set-Acl &
N/A &
Aplica listas de control de acceso desde PowerShell. &
Windows
\\
\hline

\end{longtable}


%%%%%%%%%%%%%%%%%%%%%%%%%%%%%%%%%%%%%%%%%%%%%%%%%%%%%%%%%%%%%%%%%%%%%%%%%%%%%%%

\subsubsection{ACL, SELinux y AppArmor}

Además de los permisos tradicionales, algunos sistemas utilizan mecanismos de
seguridad adicionales que pueden impedir el acceso a recursos incluso cuando
los permisos POSIX parecen correctos.

\begin{longtable}{|p{2.8cm}|p{3.2cm}|p{4cm}|p{5.4cm}|p{1.5cm}|}

\hline
Comando &
Sintaxis &
Argumentos útiles &
Particularidad &
OS
\\
\hline
\endhead

getfacl &
getfacl archivo &
-R &
Consulta las listas de control de acceso (ACL) aplicadas a un archivo o directorio. &
Linux
\\
\hline

setfacl &
setfacl &
-m, -x, -R &
Agrega, modifica o elimina ACLs. Muy utilizado cuando los permisos POSIX no son suficientes. &
Linux
\\
\hline

sestatus &
sestatus &
-v &
Consulta el estado actual de SELinux. Constituye normalmente el primer comando cuando SELinux podría estar bloqueando una aplicación. &
Linux
\\
\hline

getenforce &
getenforce &
N/A &
Indica si SELinux opera en modo Enforcing, Permissive o Disabled. &
Linux
\\
\hline

restorecon &
restorecon &
-R &
Restaura el contexto de seguridad predeterminado sobre archivos y directorios. &
Linux
\\
\hline

ls &
ls -Z &
N/A &
Muestra el contexto SELinux asociado a cada archivo. &
Linux
\\
\hline

aa-status &
aa-status &
N/A &
Consulta el estado de AppArmor y los perfiles actualmente cargados. &
Linux
\\
\hline

aa-complain &
aa-complain perfil &
N/A &
Coloca un perfil de AppArmor en modo de aprendizaje (complain). &
Linux
\\
\hline

aa-enforce &
aa-enforce perfil &
N/A &
Activa nuevamente el modo de protección de un perfil AppArmor. &
Linux
\\
\hline

\end{longtable}

%%%%%%%%%%%%%%%%%%%%%%%%%%%%%%%%%%%%%%%%%%%%%%%%%%%%%%%%%%%%%%%%%%%%%%%%%%%%%%%

\subsection{Paso 15. Transferencia y compresión de archivos}

Durante una incidencia suele ser necesario respaldar archivos, transferir logs,
copiar configuraciones entre equipos o comprimir información antes de enviarla
a otro equipo de soporte.

Las siguientes herramientas permiten realizar dichas tareas de forma segura y
eficiente.

%%%%%%%%%%%%%%%%%%%%%%%%%%%%%%%%%%%%%%%%%%%%%%%%%%%%%%%%%%%%%%%%%%%%%%%%%%%%%%%

\subsubsection{Transferencia de archivos}

\begin{longtable}{|p{2.8cm}|p{3.3cm}|p{4cm}|p{5.3cm}|p{1.5cm}|}

\hline
Comando &
Sintaxis &
Argumentos útiles &
Particularidad &
OS
\\
\hline
\endhead

scp &
scp archivo usuario@host:/ruta &
-r, -P, -i, -C &
Permite copiar archivos mediante SSH. Constituye uno de los métodos más utilizados para transferir logs entre servidores. &
Linux/macOS
\\
\hline

sftp &
sftp usuario@host &
put, get, ls, cd &
Proporciona una sesión interactiva para transferir archivos utilizando SSH. &
Linux/macOS
\\
\hline

rsync &
rsync -av origen destino &
-z, --delete, -e, --progress &
Sincroniza archivos y directorios transfiriendo únicamente las diferencias. Muy útil para respaldos y despliegues. &
Linux/macOS
\\
\hline

curl &
curl -O URL &
-o, -L, -u, --upload-file &
Además de realizar solicitudes HTTP, permite descargar y enviar archivos mediante distintos protocolos (HTTP, HTTPS, FTP, SFTP, etc.). &
Linux/macOS/Windows
\\
\hline

wget &
wget URL &
-c, -O, -q, --mirror &
Herramienta orientada principalmente a la descarga de archivos y sitios web. &
Linux
\\
\hline

ftp &
ftp servidor &
put, get &
Cliente clásico para servidores FTP. Actualmente suele reservarse para entornos heredados debido a la ausencia de cifrado. &
Linux/macOS/Windows
\\
\hline

nc &
nc host puerto &
-l, -p &
Puede utilizarse para transferencias simples de archivos entre dos equipos cuando no se dispone de otros servicios. &
Linux/macOS
\\
\hline

Copy-Item &
Copy-Item origen destino &
-Recurse &
Copia archivos y directorios mediante PowerShell. &
Windows
\\
\hline

Start-BitsTransfer &
Start-BitsTransfer URL &
Destination &
Realiza descargas utilizando el servicio BITS de Windows, permitiendo reanudación automática de transferencias. &
Windows
\\
\hline

robocopy &
robocopy origen destino &
/E, /MIR, /COPY, /MT &
Herramienta de copia robusta para Windows. Muy utilizada para respaldos y sincronización de directorios. &
Windows
\\
\hline

\end{longtable}

%%%%%%%%%%%%%%%%%%%%%%%%%%%%%%%%%%%%%%%%%%%%%%%%%%%%%%%%%%%%%%%%%%%%%%%%%%%%%%%

\subsubsection{Compresión y descompresión}

Durante el troubleshooting es frecuente comprimir archivos de log antes de
compartirlos con otros equipos o almacenarlos como evidencia.

\begin{longtable}{|p{2.8cm}|p{3.3cm}|p{4cm}|p{5.3cm}|p{1.5cm}|}

\hline
Comando &
Sintaxis &
Argumentos útiles &
Particularidad &
OS
\\
\hline
\endhead

tar &
tar -cvf respaldo.tar directorio &
-c, -x, -t, -v, -f &
Empaqueta múltiples archivos en un único archivo TAR. Normalmente se combina con herramientas de compresión como gzip o xz. &
Linux/macOS
\\
\hline

tar &
tar -xvf respaldo.tar &
-C &
Extrae el contenido de un archivo TAR. &
Linux/macOS
\\
\hline

gzip &
gzip archivo &
-d, -k &
Comprime archivos utilizando el algoritmo Gzip. &
Linux/macOS
\\
\hline

gunzip &
gunzip archivo.gz &
-c &
Descomprime archivos en formato Gzip. &
Linux/macOS
\\
\hline

bzip2 &
bzip2 archivo &
-d &
Algoritmo de compresión con mayor relación de compresión que gzip, aunque más lento. &
Linux/macOS
\\
\hline

xz &
xz archivo &
-d, -T &
Proporciona una de las mayores tasas de compresión disponibles en herramientas estándar. &
Linux/macOS
\\
\hline

zip &
zip -r respaldo.zip directorio &
-r, -9 &
Genera archivos ZIP compatibles con la mayoría de los sistemas operativos. &
Todos
\\
\hline

unzip &
unzip respaldo.zip &
-d &
Extrae archivos comprimidos en formato ZIP. &
Todos
\\
\hline

7z &
7z a respaldo.7z directorio &
x, l &
Herramienta de compresión de alta eficiencia compatible con numerosos formatos. &
Linux/Windows
\\
\hline

Compress-Archive &
Compress-Archive ruta archivo.zip &
DestinationPath &
Comprime archivos utilizando PowerShell. &
Windows
\\
\hline

Expand-Archive &
Expand-Archive archivo.zip &
DestinationPath &
Extrae archivos ZIP mediante PowerShell. &
Windows
\\
\hline

\end{longtable}

%%%%%%%%%%%%%%%%%%%%%%%%%%%%%%%%%%%%%%%%%%%%%%%%%%%%%%%%%%%%%%%%%%%%%%%%%%%%%%%

\subsubsection{Verificación de integridad}

Después de transferir archivos resulta recomendable comprobar que la
información no haya sido modificada durante el proceso.

\begin{longtable}{|p{2.8cm}|p{3.3cm}|p{4cm}|p{5.3cm}|p{1.5cm}|}

\hline
Comando &
Sintaxis &
Argumentos útiles &
Particularidad &
OS
\\
\hline
\endhead

md5sum &
md5sum archivo &
-c &
Calcula el hash MD5 de un archivo. Adecuado para verificar integridad, aunque no se recomienda para fines criptográficos. &
Linux/macOS
\\
\hline

sha1sum &
sha1sum archivo &
-c &
Calcula el hash SHA-1. Actualmente se utiliza principalmente por compatibilidad. &
Linux/macOS
\\
\hline

sha256sum &
sha256sum archivo &
-c &
Calcula el hash SHA-256. Es el algoritmo recomendado para verificar integridad durante la transferencia de archivos. &
Linux/macOS
\\
\hline

shasum &
shasum -a 256 archivo &
-a &
Versión disponible en macOS para calcular diferentes algoritmos SHA. &
macOS
\\
\hline

certutil &
certutil -hashfile archivo SHA256 &
MD5, SHA1 &
Permite calcular hashes desde la línea de comandos de Windows. &
Windows
\\
\hline

Get-FileHash &
Get-FileHash archivo &
-Algorithm &
Calcula hashes utilizando PowerShell (SHA256, SHA384, SHA512, entre otros). &
Windows
\\
\hline

\end{longtable}

%%%%%%%%%%%%%%%%%%%%%%%%%%%%%%%%%%%%%%%%%%%%%%%%%%%%%%%%%%%%%%%%%%%%%%%%%%%%%%%

\subsection{Paso 16. Verificación del entorno de ejecución}

Durante el troubleshooting es frecuente que una aplicación falle debido a una
configuración incorrecta del entorno de ejecución y no por un problema propio
del software.

Las siguientes herramientas permiten inspeccionar variables de entorno,
ejecutables disponibles, historial de comandos y otros elementos que pueden
afectar el comportamiento de una aplicación.

%%%%%%%%%%%%%%%%%%%%%%%%%%%%%%%%%%%%%%%%%%%%%%%%%%%%%%%%%%%%%%%%%%%%%%%%%%%%%%%

\subsubsection{Variables de entorno}

\begin{longtable}{|p{2.8cm}|p{3.3cm}|p{4cm}|p{5.3cm}|p{1.5cm}|}

\hline
Comando &
Sintaxis &
Argumentos útiles &
Particularidad &
OS
\\
\hline
\endhead

env &
env &
-i &
Muestra todas las variables de entorno disponibles para el proceso actual. También puede ejecutar un programa con un entorno personalizado. &
Linux/macOS
\\
\hline

printenv &
printenv &
VARIABLE &
Permite consultar el valor de una variable de entorno específica o listar todas las existentes. &
Linux/macOS
\\
\hline

export &
export VARIABLE=valor &
-n &
Define variables de entorno para la sesión actual y para los procesos hijos. &
Linux/macOS
\\
\hline

unset &
unset VARIABLE &
N/A &
Elimina variables de entorno de la sesión actual. &
Linux/macOS
\\
\hline

set &
set &
-o &
Muestra variables de shell y opciones activas. En Windows muestra las variables de entorno del sistema. &
Linux/Windows
\\
\hline

setx &
setx VARIABLE valor &
/M &
Define variables de entorno permanentes en Windows. &
Windows
\\
\hline

Get-ChildItem Env: &
Get-ChildItem Env: &
N/A &
Lista todas las variables de entorno disponibles mediante PowerShell. &
Windows
\\
\hline

\$Env:VARIABLE &
\$Env:PATH &
N/A &
Permite consultar o modificar variables de entorno desde PowerShell. &
Windows
\\
\hline

\end{longtable}

%%%%%%%%%%%%%%%%%%%%%%%%%%%%%%%%%%%%%%%%%%%%%%%%%%%%%%%%%%%%%%%%%%%%%%%%%%%%%%%

\subsubsection{Localización de ejecutables}

Una causa frecuente de incidencias consiste en ejecutar una versión diferente
del programa esperado.

\begin{longtable}{|p{2.8cm}|p{3.3cm}|p{4cm}|p{5.3cm}|p{1.5cm}|}

\hline
Comando &
Sintaxis &
Argumentos útiles &
Particularidad &
OS
\\
\hline
\endhead

which &
which python &
-a &
Muestra el ejecutable que será utilizado al invocar un comando. &
Linux/macOS
\\
\hline

whereis &
whereis nginx &
-b, -m &
Localiza ejecutables, código fuente y páginas del manual. &
Linux
\\
\hline

type &
type ls &
-a &
Indica si un comando corresponde a un ejecutable, alias, función o builtin del shell. &
Linux/macOS
\\
\hline

command &
command -v python &
-V &
Determina cómo resolverá el shell un comando determinado. &
Linux/macOS
\\
\hline

where &
where python &
N/A &
Busca ejecutables en el PATH de Windows. &
Windows
\\
\hline

Get-Command &
Get-Command git &
-All &
Obtiene información detallada sobre cmdlets, scripts y ejecutables disponibles en PowerShell. &
Windows
\\
\hline

\end{longtable}

%%%%%%%%%%%%%%%%%%%%%%%%%%%%%%%%%%%%%%%%%%%%%%%%%%%%%%%%%%%%%%%%%%%%%%%%%%%%%%%

\subsubsection{Historial y automatización}

Los siguientes comandos facilitan la reproducción de problemas y la
documentación de procedimientos.

\begin{longtable}{|p{2.8cm}|p{3.3cm}|p{4cm}|p{5.3cm}|p{1.5cm}|}

\hline
Comando &
Sintaxis &
Argumentos útiles &
Particularidad &
OS
\\
\hline
\endhead

history &
history &
-c &
Muestra el historial de comandos ejecutados por el usuario. Muy útil para reproducir acciones recientes. &
Linux/macOS
\\
\hline

fc &
fc &
-l &
Permite listar, editar y volver a ejecutar comandos del historial. &
Linux/macOS
\\
\hline

alias &
alias &
N/A &
Lista o define alias para comandos. Puede explicar comportamientos inesperados durante una sesión. &
Linux/macOS
\\
\hline

unalias &
unalias nombre &
-a &
Elimina alias previamente definidos. &
Linux/macOS
\\
\hline

tee &
comando | tee salida.log &
-a &
Guarda simultáneamente la salida de un comando en pantalla y en un archivo. Muy útil para generar evidencia durante una incidencia. &
Linux/macOS
\\
\hline

xargs &
find ... | xargs comando &
-0, -I &
Convierte la salida de un comando en argumentos para otro. Muy utilizado en comandos compuestos de troubleshooting. &
Linux/macOS
\\
\hline

nohup &
nohup comando \& &
N/A &
Permite que un proceso continúe ejecutándose aun cuando la sesión termine. &
Linux/macOS
\\
\hline

script &
script evidencia.log &
-a &
Registra toda la sesión de terminal en un archivo para documentar una incidencia. &
Linux/macOS
\\
\hline

PowerShell Start-Transcript &
Start-Transcript &
Path &
Genera un registro completo de una sesión de PowerShell. Muy útil para auditoría y documentación. &
Windows
\\
\hline

PowerShell Stop-Transcript &
Stop-Transcript &
N/A &
Finaliza la captura iniciada con Start-Transcript. &
Windows
\\
\hline

\end{longtable}

%%%%%%%%%%%%%%%%%%%%%%%%%%%%%%%%%%%%%%%%%%%%%%%%%%%%%%%%%%%%%%%%%%%%%%%%%%%%%%%

\subsubsection{Información del entorno}

\begin{longtable}{|p{2.8cm}|p{3.3cm}|p{4cm}|p{5.3cm}|p{1.5cm}|}

\hline
Comando &
Sintaxis &
Argumentos útiles &
Particularidad &
OS
\\
\hline
\endhead

pwd &
pwd &
-P &
Muestra el directorio de trabajo actual. &
Linux/macOS
\\
\hline

echo &
echo \$PATH &
N/A &
Permite visualizar variables de entorno y verificar su contenido. &
Linux/macOS
\\
\hline

dirname &
dirname ruta &
N/A &
Obtiene el directorio asociado a una ruta determinada. &
Linux/macOS
\\
\hline

basename &
basename ruta &
-s &
Extrae únicamente el nombre del archivo a partir de una ruta completa. &
Linux/macOS
\\
\hline

realpath &
realpath archivo &
N/A &
Obtiene la ruta absoluta resolviendo enlaces simbólicos. &
Linux/macOS
\\
\hline

readlink &
readlink -f archivo &
N/A &
Resuelve enlaces simbólicos y muestra el archivo real. &
Linux/macOS
\\
\hline

\end{longtable}

%%%%%%%%%%%%%%%%%%%%%%%%%%%%%%%%%%%%%%%%%%%%%%%%%%%%%%%%%%%%%%%%%%%%%%%%%%%%%%%

\subsection{Paso 17. Herramientas avanzadas de diagnóstico}

Cuando las verificaciones tradicionales no permiten identificar la causa de una
incidencia, resulta necesario utilizar herramientas que permitan inspeccionar
el comportamiento interno de procesos, binarios, bibliotecas compartidas y del
propio sistema operativo.

Estas herramientas suelen emplearse durante tareas de soporte avanzado (L2/L3),
administración de sistemas, ingeniería de confiabilidad (SRE) e ingeniería de
software.

%%%%%%%%%%%%%%%%%%%%%%%%%%%%%%%%%%%%%%%%%%%%%%%%%%%%%%%%%%%%%%%%%%%%%%%%%%%%%%%

\subsubsection{Seguimiento de llamadas al sistema}

Las siguientes herramientas permiten observar la interacción entre un proceso y
el kernel del sistema operativo.

\begin{longtable}{|p{2.8cm}|p{3.3cm}|p{4cm}|p{5.3cm}|p{1.5cm}|}

\hline
Comando &
Sintaxis &
Argumentos útiles &
Particularidad &
OS
\\
\hline
\endhead

strace &
strace programa &
-f, -p, -o, -e, -tt &
Permite observar todas las llamadas al sistema realizadas por un proceso. Es una de las herramientas más utilizadas cuando una aplicación falla sin registrar errores en sus logs.
&
Linux
\\
\hline

strace &
strace -p PID &
-f &
Se conecta a un proceso que ya se encuentra ejecutándose sin necesidad de reiniciarlo.
&
Linux
\\
\hline

ltrace &
ltrace programa &
-f, -o, -e &
Muestra las llamadas realizadas a bibliotecas compartidas (libc, OpenSSL, etc.).
&
Linux
\\
\hline

ktrace &
ktrace comando &
-i &
Equivalente disponible en algunos sistemas BSD y macOS.
&
macOS/BSD
\\
\hline

truss &
truss comando &
-p &
Herramienta equivalente presente en Solaris y algunos Unix.
&
Unix
\\
\hline

\end{longtable}

%%%%%%%%%%%%%%%%%%%%%%%%%%%%%%%%%%%%%%%%%%%%%%%%%%%%%%%%%%%%%%%%%%%%%%%%%%%%%%%

\subsubsection{Bibliotecas y ejecutables}

Cuando una aplicación no inicia, resulta habitual verificar dependencias,
bibliotecas compartidas y formato del ejecutable.

\begin{longtable}{|p{2.8cm}|p{3.3cm}|p{4cm}|p{5.3cm}|p{1.5cm}|}

\hline
Comando &
Sintaxis &
Argumentos útiles &
Particularidad &
OS
\\
\hline
\endhead

ldd &
ldd ejecutable &
-v &
Lista todas las bibliotecas dinámicas requeridas por un programa.
&
Linux
\\
\hline

file &
file ejecutable &
-i &
Determina el tipo real del archivo, arquitectura y formato.
&
Linux/macOS
\\
\hline

strings &
strings ejecutable &
-a &
Extrae cadenas ASCII presentes en archivos binarios.
&
Linux/macOS
\\
\hline

readelf &
readelf -a ejecutable &
-h, -d, -s &
Permite inspeccionar completamente un ejecutable ELF.
&
Linux
\\
\hline

objdump &
objdump -x ejecutable &
-d, -p &
Obtiene información de secciones, símbolos y código ensamblador.
&
Linux
\\
\hline

nm &
nm ejecutable &
-D &
Lista los símbolos definidos y requeridos por un ejecutable.
&
Linux/macOS
\\
\hline

otool &
otool -L ejecutable &
-hv &
Equivalente de ldd para macOS.
&
macOS
\\
\hline

dumpbin &
dumpbin /DEPENDENTS archivo.exe &
/HEADERS &
Permite inspeccionar ejecutables PE en Windows.
&
Windows
\\
\hline

\end{longtable}

%%%%%%%%%%%%%%%%%%%%%%%%%%%%%%%%%%%%%%%%%%%%%%%%%%%%%%%%%%%%%%%%%%%%%%%%%%%%%%%

\subsubsection{Depuración y análisis de procesos}

\begin{longtable}{|p{2.8cm}|p{3.3cm}|p{4cm}|p{5.3cm}|p{1.5cm}|}

\hline
Comando &
Sintaxis &
Argumentos útiles &
Particularidad &
OS
\\
\hline
\endhead

gdb &
gdb ejecutable &
-p, bt &
Permite depurar aplicaciones, inspeccionar memoria y obtener backtraces.
&
Linux
\\
\hline

gcore &
gcore PID &
N/A &
Genera un core dump de un proceso sin finalizarlo.
&
Linux
\\
\hline

pstack &
pstack PID &
N/A &
Obtiene el stack trace de un proceso en ejecución.
&
Linux
\\
\hline

coredumpctl &
coredumpctl list &
info, gdb &
Administra y consulta core dumps generados por systemd.
&
Linux
\\
\hline

lldb &
lldb ejecutable &
process attach &
Depurador principal disponible en macOS.
&
macOS
\\
\hline

WinDbg &
windbg.exe &
N/A &
Depurador avanzado para Windows.
&
Windows
\\
\hline

\end{longtable}


%%%%%%%%

%%%%%%%%%%%%%%%%%%%%%%%%%%%%%%%%%%%%%%%%%%%%%%%%%%%%%%%%%%%%%%%%%%%%%%%%%%%%%%%

\subsubsection{Análisis de rendimiento}

\begin{longtable}{|p{2.8cm}|p{3.3cm}|p{4cm}|p{5.3cm}|p{1.5cm}|}

\hline
Comando &
Sintaxis &
Argumentos útiles &
Particularidad &
OS
\\
\hline
\endhead

perf &
perf stat comando &
record, report, top &
Herramienta principal para análisis de rendimiento del kernel y aplicaciones.
&
Linux
\\
\hline

pidstat &
pidstat 1 &
-r, -u, -d &
Obtiene estadísticas detalladas por proceso.
&
Linux
\\
\hline

sar &
sar &
-u, -r, -n, -d &
Permite consultar métricas históricas de CPU, memoria, red y disco.
&
Linux
\\
\hline

dstat &
dstat &
-c, -m, -d, -n &
Presenta métricas de múltiples subsistemas simultáneamente.
&
Linux
\\
\hline

vmstat &
vmstat 1 &
N/A &
Resume continuamente el estado del sistema.
&
Linux/macOS
\\
\hline

\end{longtable}

%%%%%%%%%%%%%%%%%%%%%%%%%%%%%%%%%%%%%%%%%%%%%%%%%%%%%%%%%%%%%%%%%%%%%%%%%%%%%%%

\subsubsection{Inspección binaria}

\begin{longtable}{|p{2.8cm}|p{3.3cm}|p{4cm}|p{5.3cm}|p{1.5cm}|}

\hline
Comando &
Sintaxis &
Argumentos útiles &
Particularidad &
OS
\\
\hline
\endhead

hexdump &
hexdump -C archivo &
-n &
Visualiza archivos en formato hexadecimal.
&
Linux/macOS
\\
\hline

xxd &
xxd archivo &
-r &
Convierte archivos binarios a hexadecimal y viceversa.
&
Linux/macOS
\\
\hline

od &
od archivo &
-x, -c &
Permite visualizar archivos en distintos formatos numéricos.
&
Linux/macOS
\\
\hline

cmp &
cmp archivo1 archivo2 &
-l &
Compara dos archivos byte por byte.
&
Linux/macOS
\\
\hline

diff &
diff archivo1 archivo2 &
-u, -r &
Compara archivos de texto mostrando únicamente las diferencias.
&
Todos
\\
\hline

\end{longtable}

%%%%%%


%%%%%%%%%%%%%%%%%%%%%%%%%%%%%%%%%%%%%%%%%%%%%%%%%%%%%%%%%%%%%%%%%%%%%%%%%%%%%%%

\subsubsection{Herramientas de monitoreo avanzado}

\begin{longtable}{|p{2.8cm}|p{3.3cm}|p{4cm}|p{5.3cm}|p{1.5cm}|}

\hline
Comando &
Sintaxis &
Argumentos útiles &
Particularidad &
OS
\\
\hline
\endhead

watch &
watch comando &
-n &
Ejecuta periódicamente un comando mostrando las diferencias entre ejecuciones.
&
Linux
\\
\hline

time &
time comando &
-p &
Mide el tiempo de ejecución consumido por un programa.
&
Linux/macOS
\\
\hline

timeout &
timeout 30 comando &
-s &
Finaliza automáticamente un proceso después de un tiempo determinado.
&
Linux
\\
\hline

stdbuf &
stdbuf -oL comando &
-i, -e &
Modifica el comportamiento del buffering de entrada y salida.
&
Linux
\\
\hline

lsof &
lsof &
-p, -i, +D &
Lista todos los archivos abiertos por el sistema o por un proceso específico.
&
Linux/macOS
\\
\hline

fuser &
fuser archivo &
-k &
Determina qué proceso mantiene abierto un archivo, directorio o puerto.
&
Linux
\\
\hline

sysctl &
sysctl -a &
-w &
Consulta y modifica parámetros del kernel en tiempo de ejecución.
&
Linux/macOS
\\
\hline

ipcs &
ipcs &
-q, -m, -s &
Consulta recursos IPC (memoria compartida, colas de mensajes y semáforos).
&
Linux
\\
\hline

ipcrm &
ipcrm &
-m, -q, -s &
Elimina recursos IPC que hayan quedado huérfanos.
&
Linux
\\
\hline

\end{longtable}





\end{document}